%%%%%%%%%%%%%%%%%%%%%%%%%%%%%%%%%%%%%%%%%%%%%%%%%%%%%%%%%%%%%%%%%%%%%%
% Lớp 10 | BPT và hệ BPT bậc nhất hai ẩn | Bất phương trình bậc nhất hai ẩn
% Dạng 0D2?1-2: Miền nghiệm của BPT bậc nhất hai ẩn
% Số câu: 75  (sinh tự động bởi tach_cau.py; từ nay thầy sửa trực tiếp file này)
%%%%%%%%%%%%%%%%%%%%%%%%%%%%%%%%%%%%%%%%%%%%%%%%%%%%%%%%%%%%%%%%%%%%%%

%%%%%%%%%%%%%%%% TRẮC NGHIỆM 4 PHƯƠNG ÁN %%%%%%%%%%%%%%%%

% nguồn: MaHoaNH_TracNghiem_Lop10
\begin{ex}%[0D2N1-2]%[ma:85b4944e]
	Cho bất phương trình $2x+y\le3$. Miền nghiệm của bất phương trình trong mặt phẳng  tọa độ $Oxy$ là
	\choice {Nửa mặt phẳng chứa điểm $O$ có bờ là đường thẳng $2x+y=3$ (không kể bờ)}{Nửa mặt phẳng không chứa điểm $O$ có bờ là đường thẳng $2x+y=3$ (không kể bờ)}{\True  Nửa mặt phẳng chứa điểm $O$ có bờ là đường thẳng $2x+y=3$ (kể cả bờ)}{Nửa mặt phẳng không chứa điểm $O$ có bờ là đường thẳng $2x+y=3$ (kể cả bờ)}
	\loigiai{ 
		Vì bất phương trình có dấu \lq\lq$\le$\rq\rq~ nên miền nghiệm \textbf{kể cả  bờ}.\\
		Thay $x=0$; $y=0$ vào bất  phương trình $2(0)+0 \le 3 \Leftrightarrow 0 \le 3$ (đúng).\\
		Vậy miền nghiệm  \textbf{chứa điểm $O$}.
	 }
\end{ex}

% nguồn: MaHoaNH_TracNghiem_Lop10
\begin{ex}%[0D2N1-2]%[ma:307a4de7]
	Cho bất phương  trình $2x-y\le3$. Miền nghiệm của bất phương trình trong mặt phẳng tọa độ $Oxy$ là
	\choice {Nửa mặt phẳng chứa điểm $O$ có bờ là đường thẳng $2x-y=3$ (không kể bờ)}{Nửa mặt phẳng không chứa điểm $O$ có bờ là đường thẳng $2x-y=3$ (không kể bờ)}{\True Nửa mặt phẳng chứa điểm $O$ có bờ là đường thẳng $2x-y=3$ (kể cả bờ)}{Nửa mặt phẳng không chứa điểm $O$ có bờ là đường thẳng $2x-y=3$  (kể cả bờ)}
	\loigiai{ 
		Vì bất phương trình có dấu "$\le$"  nên miền nghiệm \textbf{kể cả bờ} là đường thẳng $2x-y=3$.\\
		 Thay $x=0;$ $y=0$ vào bất phương trình: $2(0)-0 \le 3 \Leftrightarrow 0 \le 3$ (đúng).\\
		Do đó, miền nghiệm \textbf{chứa điểm  $O(0;0)$}.\\
		Kết hợp lại, miền nghiệm là nửa mặt  phẳng chứa điểm $O$, kể cả bờ là đường thẳng $2x-y=3$.
	 }
\end{ex}

% nguồn: MaHoaNH_TracNghiem_Lop10
\begin{ex}%[0D2N1-2]%[ma:f97bb356]
	Trong các cặp số sau đây, cặp nào \textbf{không} là nghiệm của bất phương trình $2x+y<1$? 
	\choice 
		 {$(-2;-5)$}
		{$(3;-7)$} 
		{\True $(0;1)$} 
		{$(0;-6)$} 
	 \loigiai{
		 Chỉ có cặp số $(0;1)$ không thỏa bất phương trình.
	} 
\end{ex}

% nguồn: MaHoaNH_TracNghiem_Lop10
\begin{ex}%[0D2N1-2]%[ma:fc0e7017]
	Bất phương trình $2x^2-5x+2\geq 0$ có tập  nghiệm là
	 \choice
		 {$\left(-\infty;\dfrac{1}{2} \right]$}
		{$\left(-\infty;\dfrac{1}{2}\right)\cup  \left(2;+\infty \right)$}
		{\True  $\left(-\infty;\dfrac{1}{2}\right]\cup \left[2;+\infty \right)$}
		{$\left[\dfrac{1}{2};2\right]$} 
	 \loigiai{
		Ta có $2x^2-5x+2\geq 0\Leftrightarrow\hoac{&x\geq2\\&x\leq \dfrac{1}{2}.}$\\ Vậy tập nghiệm của phương trình là $S=\left(-\infty;\dfrac{1}{2}\right]\cup \left[2;+\infty  \right)$.
	 }
\end{ex}

% nguồn: MaHoaNH_TracNghiem_Lop10
\begin{ex}%[0D2N1-2]%[ma:f4857208]
	Miền nghiệm  của bất phương trình $-3x+y+1\le 0$ không chứa điểm nào sau đây?
	 \choice
		{\True $A(1;2)$} 
		 {$B(2;1)$}
		 {$D(3;1)$}
		{$C\left(1;\dfrac{1}{2}\right)$} 
	\loigiai{ 
		Vì $-3\cdot 1+2+2 > 0$ nên miền nghiệm của bất phương trình $-3x+y+1\le  0$ không chứa điểm $A(1;2)$.
	} 
\end{ex}

% nguồn: MaHoaNH_TracNghiem_Lop10
\begin{ex}%[0D2N1-2]%[ma:ed346c4d]
	Bất phương trình $-x^2+2x+3<0$ có tập  nghiệm là
	\choice 
		 {$(-1;3)$}
		 {$(-3;1)$}
		{$[-1;3]$} 
		{\True$(-\infty;-1) \cup  (3;+\infty)$}
	\loigiai{ 
		Ta có $-x^2+2x+3<0\Leftrightarrow \hoac{&x<-1\\&x>3.}$ \\ 
		Vậy tập nghiệm của bất phương trình là $S=(-\infty;-1)\cup  (3;+\infty)$.
	} 
\end{ex}

% nguồn: MaHoaNH_TracNghiem_Lop10
\begin{ex}%[0D2N1-2]%[ma:3bbecf84]
	Điểm nào sau đây thuộc miền nghiệm  của bất phương trình $2x+y-3>0$?
	 \choice
		{$Q\left(-1;-3\right)$} 
		{\True $M\left(1;\dfrac{3}{2}\right)$} 
		 {$N\left(1;1\right)$}
		{$P\left(-1;\dfrac{3}{2}\right)$} 
	\loigiai{ 
		\begin{itemize}
			\item Xét điểm $Q\left(-1;-3\right)$, ta  có $2\cdot(-1)+(-3)-3=-7<0$. Do đó điểm $Q$ không thuộc miền nghiệm của bất phương trình $2x+y-3>0$.
			\item Xét điểm $M\left(1;\dfrac{3}{2}\right)$, ta  có $2\cdot1+\dfrac{3}{2}-3=\dfrac{1}{2}>0$. Do đó điểm $M$ thuộc miền nghiệm của bất phương trình $2x+y-3>0$.
			\item Xét điểm $N\left(1;1\right)$, ta có $2\cdot1+1-3=0$.  Do đó điểm $N$ không thuộc miền nghiệm của bất phương trình $2x+y-3>0$.
			\item Xét điểm $P\left(-1;\dfrac{3}{2}\right)$, ta có $2\cdot(-1)+\dfrac{3}{2}-3=-\dfrac{7}{2}<0$. Do đó điểm  $P$ không thuộc miền nghiệm của bất phương trình $2x+y-3>0$.
		\end{itemize}
	} 
\end{ex}

% nguồn: MaHoaNH_TracNghiem_Lop10
\begin{ex}%[0D2N1-2]%[ma:f7c70192]
	Cặp số nào dưới đây là nghiệm của bất  phương trình $2x-y>3$?
	 \choice
		{\True  $(3;1)$}
		{$(-2;1)$} 
		{$(1;1)$} 
		 {$(2;1)$}
	\loigiai{ 
		\begin{itemize}
			\item Xét cặp số $(3;1)$, ta có $2\cdot 3 -1 = 5 >3\Rightarrow (3;1)$ là nghiệm của bất  phương trình $2x-y>3$.
			\item Xét cặp  số $(-2;1)$, ta có $2\cdot (-2) -1 =-5<3\Rightarrow (-2;1)$ không là nghiệm của bất phương trình $2x-y>3$.
			\item Xét cặp số $(1; 1)$, ta có $2\cdot 1 -1 =1<3\Rightarrow (1; 1)$ không là nghiệm của bất phương trình  $2x-y>3$.
			\item Xét cặp số $(2;1)$, ta có $2\cdot 2 -1 =3\Rightarrow (2;1)$ không là một nghiệm của bất phương trình $2x-y>3$. 
		\end{itemize}
		Vậy cặp số $(3; 1)$ là nghiệm của bất phương  trình $2 x- y <3$.
	 }
\end{ex}

% nguồn: MaHoaNH_TracNghiem_Lop10
\begin{ex}%[0D2N1-2]%[ma:01cf5ce9]
	Cặp số nào sau đây \textbf{không} là nghiệm của bất phương trình  $x-2y\ge 5$?
	\choice 
		 {$(3;-1)$}
		{\True $(-1;4)$} 
		{$(2;-3)$} 
		 {$(1;-2)$}
	\loigiai{ 
		\begin{itemize}
			\item Xét cặp số $(3;-1)$, ta  có $3 -2\cdot(-1) = 5\Rightarrow (3;-1)$ là nghiệm của bất phương trình $x-2y\ge 5$.
			\item Xét cặp số $(-1;4)$, ta có $-1 -2\cdot4 =-9<5\Rightarrow (-1;4)$ không là  nghiệm của bất phương trình $x-2y\ge 5$.
			\item Xét cặp số $(2; -3)$, ta có $2 -2\cdot(-3) =8>5\Rightarrow (2; -3)$ là nghiệm của bất phương  trình $x-2y\ge 5$.
			\item Xét cặp số $(1; -2)$, ta có $1 -2\cdot(-2) =5\Rightarrow (1; -2)$ là một  nghiệm của bất phương trình $x-2y\ge 5$.
		\end{itemize}
		Vậy cặp số  $(-1; 4)$ là không nghiệm của bất phương trình $x-2y\ge 5$.
	 }
\end{ex}

% nguồn: MaHoaNH_TracNghiem_Lop10
\begin{ex}%[0D2N1-2]%[ma:76a2a247]
	Điểm $A(1;-3)$ là điểm thuộc miền nghiệm của bất phương trình  nào sau đây?
	 \choice
		{$-3x+2y-3>0$} 
		 {$3x-y\le 0$}
		{\True $3x-y>0$} 
		 {$y-2x>-4$}
	\loigiai{ 
		\begin{itemize}
			\item Xét bất phương trình $-3x+2y-3>0$, ta có $-3\cdot1 + 2\cdot (-3)-3=-12<0$. Do đó điểm $A(1; -3)$ không thuộc miền nghiệm của bất  phương trình $-3x+2y-3>0$.
			\item Xét bất phương trình $3x-y\le 0$, ta có $3\cdot1 -(-3)=6>0$. Do đó điểm $A(1;  -3)$ không thuộc miền nghiệm của bất phương trình $3x-y\le 0$.
			\item Xét bất phương trình $3x-y>0$, ta có $3\cdot1 -(-3)=6>0$. Do đó điểm $A(1; -3)$  thuộc miền nghiệm của bất phương trình $3x-y>0$.
			\item Xét bất phương trình $y-2x>-4$, ta có $-3 -  2\cdot 1=-5<-4$. Do đó điểm $A(1; -3)$ không thuộc miền nghiệm của bất phương trình $y-2x>-4$.
		\end{itemize}
		Vậy điểm $A(1; -3)$ thuộc miền nghiệm của bất phương trình  $3x-y>0$.
	} 
\end{ex}

% nguồn: MaHoaNH_TracNghiem_Lop10
\begin{ex}%[0D2N1-2]%[ma:60538c23]
	Miền nghiệm của bất phương trình $5(x+2)-9<2x-2y+7$ \textbf{không}  chứa điểm nào trong các điểm sau?
	\choice 
		 {\True $(2;3)$}
		 {$(-2;1)$}
		{$(2;-1)$} 
		{$(0;0)$} 
	 \loigiai{
		Ta có $5(x+2)-9<2x-2y+7\Leftrightarrow 5x+10-9<2x-2y+7\Leftrightarrow 3x+2y<6$. 
		\begin{itemize}
			\item Xét  điểm $(2;3)$, ta có $3\cdot2 +2\cdot3=12>6$. Do đó miển nghiệm của bất phương trình đã cho không chứa điểm $(2; 3)$.
			\item  Xét điểm $(-2;1)$, ta có $3\cdot(-2) +2\cdot1=-4<6$. Do đó miển nghiệm của bất phương trình đã cho chứa điểm $(-2;1)$.
			\item Xét điểm $(2;-1)$, ta có $3\cdot2 +2\cdot(-1)=4>6$. Do đó miển nghiệm của  bất phương trình đã cho chứa điểm $(2;-1)$.
			\item Xét điểm $(0;0)$, ta  có $3\cdot0 +2\cdot0=0<6$. Do đó miển nghiệm của bất phương trình đã cho chứa điểm $(0;0)$.
		\end{itemize}
		Vậy miền nghiệm của bất phương trình  $5(x+2)-9<2x-2y+7$ không chứa điểm $(2; 3)$.
	} 
\end{ex}

% nguồn: MaHoaNH_TracNghiem_Lop10
\begin{ex}%[0D2N1-2]%[ma:f2870407]
	 Miền nghiệm của bất phương trình $3(x-1)+4(y-2)<5x-3$ là nửa mặt phẳng chứa điểm nào sau đây?
	\choice 
		{$Q(-5;3)$} 
		{\True  $O(0;0)$}
		 {$N(-4;2)$}
		 {$P(-2;2)$}
	 \loigiai{
		Ta  có $3(x-1)+4(y-2)<5x-3\Leftrightarrow -x+2y-4<0$.
		\begin{itemize}
			\item Xét điểm $Q(-5;3)$, ta có $-(-5)+2\cdot3-4=7>0$. Do  đó miền nghiệm của bất phương trình đã cho là nửa mặt phẳng không chứa điểm $Q(-5;3)$.
			\item Xét điểm $O(0;0)$, ta có $-0+2\cdot0-4=-4<0$ . Do  đó miền nghiệm của bất phương trình đã cho là nửa mặt phẳng chứa điểm $O(0;0)$.
			\item Xét điểm $N(-4;2)$,  ta có $-(-4)+2\cdot2-4=4>0$. Do đó miền nghiệm của bất phương trình đã cho là nửa mặt phẳng không chứa điểm $N(-4;2)$.
			\item Xét điểm $P(-2;2)$, ta có $-(-2)+2\cdot2-4=2>0$. Do đó miền nghiệm của bất phương trình  đã cho là nửa mặt phẳng không chứa điểm $P(-2;2)$.
		\end{itemize}
		Vậy miền nghiệm của bất phương trình đã cho là nửa mặt  phẳng chứa điểm $O(0;0)$.
	} 
\end{ex}

% nguồn: MaHoaNH_TracNghiem_Lop10
\begin{ex}%[0D2N1-2]%[ma:cc8b88f0]
	Điểm $A\left(-1;3\right)$ là điểm {\bf không thuộc} miền nghiệm của bất phương trình nào sau đây? 
	\choice 
		 {\True $-3x+2y-4<0$}
		{$x+3y>0$} 
		{$3x-y\le 0$} 
		 {$2x-y+4<0$}
	\loigiai{ 
		Thay lần lượt tọa độ $A\left(-1;3\right)$ vào bất phương trình ở đáp  án, ta có
		\begin{itemize}
			\item Xét  bất phương trình $-3x+2y-4<0$, ta có $-3\cdot(-1)+2\cdot3-4=5>0$.
			\item Xét bất phương trình $ x+3y>0$, ta có $-1+3\cdot3-4=4>0$. 
			\item Xét bất  phương trình $3x-y\le 0$, ta có $3\cdot(-1)-3-4=-10\le 0$.
			\item Xét bất phương trình  $2x-y+4<0$, ta có $2\cdot(-1)-3-4=-9<0$.
		\end{itemize}
		Vậy điểm $A\left(-1;3\right)$ là  điểm không thuộc miền nghiệm của bất phương trình $-3x+2y-4<0$.
	 }
\end{ex}

% nguồn: MaHoaNH_TracNghiem_Lop10
\begin{ex}%[0D2N1-2]%[ma:798979a9]
	Miền nghiệm của bất phương trình  $\left(1+\sqrt{3}\right)x-\left(1-\sqrt{3}\right)y\ge 2$ chứa điểm nào sau đây?
	 \choice
		 {\True $A\left(1; -1\right)$}
		{$B\left(-1;  -1\right)$}
		{$C\left(-1; ,1\right)$} 
		{$D\left(-\sqrt{3}; \sqrt{3}\right)$} 
	\loigiai{ 
		\begin{itemize}
			 \item Xét điểm $A\left(1; -1\right)$, ta có $\left(1+\sqrt{3}\right)\cdot1-\left(1-\sqrt{3}\right)\cdot1=2\sqrt{3}>2$.\\
			 Do đó miền nghiệm của bất phương trình $\left(1+\sqrt{3}\right)x-\left(1-\sqrt{3}\right)y\ge 2$ chứa điểm $A$.
			\item Xét điểm $B\left(-1; -1\right)$, ta có $\left(1+\sqrt{3}\right)\cdot(-1)-\left(1-\sqrt{3}\right)\cdot(-1)=-2\sqrt{3}<2$.\\ Do đó miền nghiệm của bất phương trình $\left(1+\sqrt{3}\right)x-\left(1-\sqrt{3}\right)y\ge 2$ không chứa  điểm $B$.
			\item  Xét điểm $C\left(-1; ,1\right)$, ta có $\left(1+\sqrt{3}\right)\cdot(-1)-\left(1-\sqrt{3}\right)\cdot1=-2<2$.\\
			Do đó miền nghiệm của bất phương trình $\left(1+\sqrt{3}\right)x-\left(1-\sqrt{3}\right)y\ge 2$ không chứa  điểm $C$.
			\item Xét điểm $D\left(-\sqrt{3};  \sqrt{3}\right)$, ta có $\left(1+\sqrt{3}\right)\cdot(-\sqrt{3})-\left(1-\sqrt{3}\right)\cdot\sqrt{3}=-2\sqrt{3}<2$.\\
			Do đó miền nghiệm của  bất phương trình $\left(1+\sqrt{3}\right)x-\left(1-\sqrt{3}\right)y\ge 2$ không chứa điểm $D$.
		\end{itemize}
	 }
\end{ex}

% nguồn: MaHoaNH_TracNghiem_Lop10
\begin{ex}%[0D2N1-2]%[ma:68c505ab]
	Miền nghiệm của bất phương trình $3x+2\left(y+3\right)>4\left(x+1\right)-y+3$ là  phần mặt phẳng chứa điểm nào?
	 \choice
		{$\left(3;0\right)$} 
		{$\left(3;1\right)$} 
		{\True  $\left(1;1\right)$}
		 {$\left(0;0\right)$}
	\loigiai{ 
		Ta có 
		\begin{eqnarray*}
			&&  3x+2\left(y+3\right)>4\left(x+1\right)-y+3\\
			 &\Leftrightarrow& 3x +2y+6> 4x+4-y+3\\
			 &\Leftrightarrow& -x +3y-1>0.
		\end{eqnarray*}
		\begin{itemize}
			\item Xét  điểm $\left(3;0\right)$, ta có $-3+3\cdot0-1=-4<0$. Do đó miền nghiệm của bất phương trình $-x +3y-1>0$ không chứa điểm $\left(3;0\right)$.
			\item Xét điểm $\left(3;1\right)$, ta có $-3+3\cdot1-1=-1<0$. Do đó miền  nghiệm của bất phương trình $-x +3y-1>0$ không chứa điểm $\left(3;1\right)$.
			 \item Xét điểm $\left(1;1\right)$, ta có $-1+3\cdot1-1=1>0$. Do đó miền nghiệm của bất phương trình $-x +3y-1>0$ chứa điểm $\left(1;1\right)$.
			\item Xét điểm $\left(0;0\right)$, ta có $-0+3\cdot0-1=-1<0$. Do đó miền nghiệm của bất phương trình $-x +3y-1>0$ không chứa điểm $\left(0;0\right)$. 
		\end{itemize}
	 }
\end{ex}

% nguồn: MaHoaNH_TracNghiem_Lop10
\begin{ex}%[0D2N1-2]%[ma:fe8d6d40]
	Miền nghiệm của bất phương trình $x-2+2\left(y-1\right)>2x+4$  chứa điểm nào sau đây?
	 \choice
		 {$A\left(1; 1\right)$}
		{\True $B\left(1;  5\right)$}
		{$C\left(4; 3\right)$} 
		{$D\left(0; 4\right)$} 
	 \loigiai{
		 Ta có
		\begin{eqnarray*}
			& & x-2+2\left(y-1\right)>2x+4\\ 
			 &\Leftrightarrow& x-2+2y-2>2x+4\\
			&\Leftrightarrow&  -x + 2y>8.
		\end{eqnarray*}
		\begin{itemize}
			\item Xét điểm $A\left(1; 1\right)$, ta có $-1 + 2\cdot1=1<8$. Do đó miền  nghiệm của bất phương trình $-x + 2y>8$ không chứa điểm $A$.
			\item Xét điểm $B\left(1; 5\right)$, ta có $1 + 2\cdot5=11>8$. Do đó miền nghiệm của bất phương trình $-x +  2y>8$ chứa điểm $B$.
			\item Xét điểm $C\left(4; 3\right)$, ta có $-4  + 2\cdot3=2<8$. Do đó miền nghiệm của bất phương trình $-x + 2y>8$ không chứa điểm $C$.
			\item Xét điểm $D\left(0; 4\right)$, ta có $-0 + 2\cdot4=8$. Do đó miền nghiệm của bất phương trình $-x + 2y>8$ chứa điểm $D$. 
		\end{itemize}
	 }
\end{ex}

% nguồn: MaHoaNH_TracNghiem_Lop10
\begin{ex}%[0D2N1-2]%[ma:b01c427a]
	Miền nghiệm  của bất phương trình $5\left(x+2\right)-9<2x-2y+7$ là phần mặt phẳng {\bf không chứa} điểm nào?
	 \choice
		 {$\left(-2;1\right)$}
		 {\True $\left(2;3\right)$}
		{$\left(2;-1\right)$} 
		{$\left(0;0\right)$} 
	\loigiai{ 
		Ta có 
		\begin{eqnarray*}
			& &  5\left(x+2\right)-9<2x-2y+7\\
			&\Leftrightarrow& 5x +10 -9<  2x-2y+7\\
			 &\Leftrightarrow& 3x + 2y-6<0.
		\end{eqnarray*}
		\begin{itemize}
			\item Xét điểm $\left(-2;1\right)$, ta có $3\cdot(-2)+2\cdot1-6=-10<0$. Do đó miền nghiệm của bất phương trình $3x + 2y-6<0$ chứa  điểm $\left(-2;1\right)$.
			\item Xét điểm $\left(2;3\right)$, ta có $3\cdot2+2\cdot3-6=6>0$. Do đó miền nghiệm của bất phương trình $3x + 2y-6<0$ không chứa  điểm $\left(2;3\right)$.
			\item Xét điểm $\left(2;-1\right)$, ta có $3\cdot2+2\cdot(-1)-6=-2<0$. Do đó miền nghiệm của bất phương trình $3x  + 2y-6<0$ chứa điểm $\left(2;-1\right)$.
			\item Xét điểm $\left(0;0\right)$,  ta có $3\cdot0+2\cdot0-6=-6<0$. Do đó miền nghiệm của bất phương trình $3x + 2y-6<0$ chứa điểm $\left(0;0\right)$.
		\end{itemize}
	 }
\end{ex}

% nguồn: MaHoaNH_TracNghiem_Lop10
\begin{ex}%[0D2N1-2]%[ma:edc2ef6b]
	Phần  không bị gạch sọc trong hình vẽ bên dưới là miền nghiệm của bất phương trình nào?
	\begin{center}
		\begin{tikzpicture}[scale=.8,>=stealth,font=\footnotesize,line join=round,line cap=round]
				 \def\a{3/2} \def\b{3}
				\def\xmin{-4} \def\xmax{2} 
				\def\ymin{-1}  \def\ymax{4}
				\draw[->]  (\xmin,0)--(\xmax+.2,0) node [below]{$x$};
				 \draw[->] (0,\ymin)--(0,\ymax+.2) node [left]{$y$};
				\fill (0,0) circle (1pt) node[below left]{$O$}; 
				\clip (\xmin,\ymin) rectangle (\xmax,\ymax); 
				\fill[pattern=north west  lines]plot (\x,{\a*(\x)+\b})--(\xmin,\ymax)--(\xmin,\ymin)--cycle;
				 \draw[smooth,samples=300] plot (\x,{\a*(\x)+\b});
				\draw[fill=black] (0,3)  circle (1pt) node [right]{$3$};
				\draw[fill=black]  (-2,0) circle (1pt) node [below]{$-2$};
			\end{tikzpicture}
	\end{center}
	\choice 
		{$ 3x-2y<-6$} 
		{\True  $ 3x-2y>-6$}
		{$  3x-2y>0$}
		 {$ 3x-2y<0$}
	 \loigiai{
		Đường  thẳng $\Delta\colon ax+by=c$ đi qua hai điểm $\left(-2;0\right)$, $\left(0;3\right)$.\\
		Suy ra $\heva{&a\cdot(-2)+b\cdot0=c\\&a\cdot0+b\cdot3=c}\Rightarrow -2a=3b$.\\ 
		Chọn $\heva{&a=3\\&b=-2}\Rightarrow c=-6\Rightarrow \Delta\colon 3x-2y=-6$.\\ 
		Ta thấy điểm $\left(0;0\right)$ thuộc  miền nghiệm của bất phương trình (phần không gạch sọc).\\
		Do đó không phần  gạch sọc là miền nghiệm của bất phương trình $3x-2y>-6$.
	 }
\end{ex}

% nguồn: MaHoaNH_TracNghiem_Lop10
\begin{ex}%[0D2N1-2]%[ma:055195ee]
	Phần không bị gạch sọc trong hình  vẽ bên dưới là miền nghiệm của bất phương trình nào?
	\begin{center}
		\begin{tikzpicture}[scale=.8,>=stealth,font=\footnotesize,line join=round,line cap=round]
				\def\a{1/2} \def\b{-2} 
				\def\xmin{-2} \def\xmax{5} 
				 \def\ymin{-3} \def\ymax{2}
				\draw[->] (\xmin,0)--(\xmax+.2,0) node  [below]{$x$};
				\draw[->]  (0,\ymin)--(0,\ymax+.2) node [left]{$y$};
				\fill (0,0)  circle (1pt) node[below left]{$O$};
				\clip  (\xmin,\ymin) rectangle (\xmax,\ymax);
				\fill[pattern=north west  lines] (\xmax,\ymax)--(\xmin,\ymax)--plot (\x,{\a*(\x)+\b})--(5,.5)--cycle;
				\draw[smooth,samples=300] plot (\x,{\a*(\x)+\b}); 
				\draw[fill=black] (0,-2) circle (1pt) node  [below right]{$-2$};
				\draw[fill=black] (4,0) circle (1pt)  node [below]{$4$};
			\end{tikzpicture}
	\end{center}
	 \choice
		{$2x-4y<8$} 
		 {\True $2x-4y>8$}
		{$2x-4y>-5$} 
		 {$2x-4y>-3$}
	\loigiai{ 
		Đường thẳng $\Delta\colon ax+by=c$ đi qua hai điểm $\left(4;0\right)$,  $\left(0;-2\right)$.\\
		Suy  ra $\heva{&a\cdot4+b\cdot0=c\\&a\cdot0+b\cdot(-2)=c}\Rightarrow 4a=-2b$.\\
		Chọn  $\heva{&a=2\\&b=-4}\Rightarrow c=8\Rightarrow \Delta\colon 2x-8y=8$.\\
		Ta  thấy điểm $\left(0;0\right)$ không thuộc miền nghiệm của bất phương trình (phần gạch sọc).\\
		Do đó không phần gạch sọc là miền nghiệm của bất phương trình $2x-8y>8$. 
	 }
\end{ex}

% nguồn: MaHoaNH_TracNghiem_Lop10
\begin{ex}%[0D2N1-2]%[ma:9019dca8]
	Phần bị gạch sọc trong hình vẽ bên dưới là miền nghiệm của bất phương trình  nào?\\
	\begin{center}
		\begin{tikzpicture}[scale=.8,>=stealth,font=\footnotesize,line join=round,line cap=round]
				 \def\a{2} \def\b{-3}
				\def\xmin{-2} \def\xmax{4} 
				\def\ymin{-4}  \def\ymax{2}
				 \draw[->] (\xmin,0)--(\xmax+.2,0) node [below]{$x$};
				\draw[->] (0,\ymin)--(0,\ymax+.2) node  [left]{$y$};
				\fill (0,0) circle (1pt)  node[below left]{$O$};
				\clip (\xmin,\ymin) rectangle (\xmax,\ymax); 
				\fill[pattern=north  west lines]plot (\x,{\a*(\x)+\b})--(\xmax,\ymax)--(\xmax,\ymin)--cycle;
				 \draw[smooth,samples=300] plot (\x,{\a*(\x)+\b});
				\draw[fill=black] (0,-3) circle (1pt) node [left]{$-3$}; 
				\draw[fill=black] (1.5,0)  circle (1pt) node [above]{$\dfrac{3}{2}$};
			\end{tikzpicture}
	\end{center}
	\choice 
		 {\True $ 2x-y>3$}
		 {$ 2x-y<3$}
		{$  x-2y<3$}
		{$  x-2y>3$}
	 \loigiai{
		Đường thẳng $\Delta\colon ax+b=c$ đi qua hai điểm $\left(\dfrac{3}{2};0\right)$, $\left(0;-3\right)$.\\ 
		Suy ra $\heva{&a\cdot\dfrac{3}{2}+b\cdot0=c\\&a\cdot0+b\cdot(-3)=c}\Rightarrow \dfrac{3}{2}a=-3b$.\\ 
		Chọn  $\heva{&a=2\\&b=-1}\Rightarrow c=3\Rightarrow \Delta\colon 2x-y=3$.\\
		Ta thấy điểm $\left(0;0\right)$ không thuộc miền nghiệm của bất phương trình (phần gạch sọc).\\ 
		Do đó  phần gạch sọc trong hình vẽ là miền nghiệm của bất phương trình $2x-y>3$.
	 }
\end{ex}

% nguồn: MaHoaNH_TracNghiem_Lop10
\begin{ex}%[0D2H1-2]%[ma:94edf2a1]
	 Miền nghiệm của bất phương trình $-x+2+2(y-2) < 2(1-x)$ là nửa mặt phẳng \textbf{không} chứa điểm nào trong các điểm sau?
	\choice 
		{$A\left(0;0\right)$} 
		 {$B\left(1;1\right)$}
		 {\True $C\left(4;2\right)$}
		 {$D\left(1;-1\right)$}
	\loigiai{Ta  có $-x+2(y-2)<2(1-x) \Leftrightarrow -x+2 +2y-4 -2+2x<0 \Leftrightarrow x+2y-4<0$.\\
	Ta có $4+2 \cdot 2-4>0$ nên miền nghiệm của bất phương trình $x+2y-4<0$  không chứa điểm $C\left(4;2\right)$.}
\end{ex}

% nguồn: MaHoaNH_TracNghiem_Lop10
\begin{ex}%[0D2H1-2]%[ma:37dad15e]
	Cho hai bất phương trình $2x+y<3\quad (1)$ và $-x+3y>5\quad (2)$ và điểm  $M(0;1)$. Kết luận nào sau đây \textbf{đúng}?
	\choice 
		{Điểm $M$ thuộc miền  nghiệm của bất phương trình $(1)$ và $(2)$}
		{\True Điểm $M$ thuộc miền nghiệm của bất phương trình $(1)$  nhưng không thuộc miền nghiệm của bất phương trình $(2)$}
		{Điểm $M$ không thuộc  miền nghiệm của bất phương trình $(1)$ nhưng thuộc miền nghiệm của bất phương trình $(2)$}
		 {Điểm $M$ không thuộc miền nghiệm của cả hai bất phương trình $(1)$ và $(2)$}
	\loigiai{ 
		Thay điểm $M(0;1)$ vào bất phương trình  $(1)$ ta được $2\cdot 0+1<3\Leftrightarrow 1<3$ (đúng).\\
		Thay điểm $M(0;1)$ vào bất phương trình $(2)$ ta được $ 0+3\cdot  1>5\Leftrightarrow 3>5$ (vô lý).
	 }
\end{ex}

% nguồn: MaHoaNH_TracNghiem_Lop10
\begin{ex}%[0D2H1-2]%[ma:a5b9efe0]
	Cho bất  phương trình $x+2y\le 3$. Khẳng định nào sau đây là đúng?
	 \choice
		{\True Miền nghiệm của bất phương trình là nửa mặt phẳng bờ $d:x+2y=3$ chứa  gốc tọa độ}
		{Miền nghiệm của bất phương trình là nửa mặt phẳng bờ $d:x+2y=3$ không chứa gốc  tọa độ}
		{Miền nghiệm của bất phương trình là nửa mặt  phẳng bờ $d:x+2y=-3$ chứa gốc tọa độ}
		 {Miền nghiệm của bất phương trình là nửa mặt phẳng bờ $d:x+2y=-3$ không chứa gốc tọa độ}
	\loigiai{ 
		Thay $(0;0)$ vào bất phương trình $x+2y\le 3\Rightarrow 0 +2\cdot 0\le  3\Leftrightarrow 0 \le 3$ (đúng).\\
		Vậy miền  nghiệm chứa điểm $(0;0)$.
	 }
\end{ex}

% nguồn: MaHoaNH_TracNghiem_Lop10
\begin{ex}%[0D2H1-2]%[ma:f14b5a0a]
	Miền nghiệm của bất phương trình $2x-3y>5$ là nửa mặt phẳng (không kể đường thẳng $d:2x-3y=5$) \textbf{không} chứa điểm  có tọa độ nào sau đây?
	 \choice
		{\True $(0;0)$} 
		 {$(3;0)$}
		{$(1;-2)$} 
		 {$(-3;-4)$}
	\loigiai{ 
		\begin{itemize}
			\item Xét điểm $(0;0)$, ta có $2\cdot0 -3\cdot0=0<5$. Do đó miền nghiệm của bất phương  trình đã cho không chứa điểm $(0; 0)$.
			\item  Xét điểm $(3;0)$, ta có $2\cdot3 -3\cdot0=6>5$. Do đó miền nghiệm của bất phương trình đã cho chứa điểm $(3;0)$.
			\item Xét điểm $(1;-2)$, ta có  $2\cdot1 -3\cdot(-2)=8>5$. Do đó miền nghiệm của bất phương trình đã cho chứa điểm $(1;-2)$.
			\item Xét điểm $(-3;-4)$, ta có $2\cdot(-3) -3\cdot(-4)=6>5$. Do đó miền nghiệm của  bất phương trình đã cho chứa điểm $(-3;-4)$.
		\end{itemize}
		Vậy miền nghiệm của bất phương trình $2x-3y>5$ không chứa điểm  $(0; 0)$.
	 }
\end{ex}

% nguồn: MaHoaNH_TracNghiem_Lop10
\begin{ex}%[0D2H1-2]%[ma:c36fc3be]
	Khẳng định nào dưới đây là khẳng định  \textbf{đúng}?
	 \choice
		{Bất phương trình $3x-2 y>0$ có miền nghiệm là nửa mặt phẳng  không kể bờ là đường thẳng $3x-2 y=0$ chứa điểm $O$}
		{Bất phương trình $2 x+y>1$ có miền nghiệm là nửa mặt phẳng không kể bờ là đường  thẳng $2x+ y =1$ chứa điểm $O$}
		{Bất phương trình $-2 x+y+1 \le 0$ có miền nghiệm là nửa mặt phẳng kể cả bờ là đường thẳng  $-2 x+y+1=0$ chứa điểm $O$}
		{\True Bất  phương trình $2 x-3 y+5 \ge 0$ có miền nghiệm là nửa mặt phẳng kể cả bờ là đường thẳng $2 x-3 y+5=0$ chứa điểm $O$}
	\loigiai{ 
		\begin{itemize}
			\item Xét bất phương trình $3x-2 y>0$, ta có $3\cdot0 -2\cdot 0 = 0$. Do đó miền  nghiệm của bất phương trình $3x-2 y>0$ không chứa điểm $O(0;0)$.
			\item Xét bất phương trình $2 x+y>1$, ta có $2\cdot0 +0= 0<1$. Do đó  miền nghiệm của bất phương trình $2 x+y>1$ không chứa điểm $O(0;0)$.
			\item Xét bất phương trình $-2 x+y+1 \le 0$, ta có $-2\cdot0 +0+1= 1>0$. Do đó miền nghiệm của bất phương trình $-2 x+y+1 \le 0$ không  chứa điểm $O(0;0)$.
			\item  Xét bất phương trình $2 x-3 y+5 \ge 0$, ta có $2\cdot0 -3\cdot0+5= 5>0$. Do đó miền nghiệm của bất phương trình $2 x-3 y+5 \ge 0$ chứa điểm $O(0;0)$.
		\end{itemize}
		Vậy bất phương trình $-2 x+y+1 \le 0$ có miền nghiệm là nửa mặt phẳng kể cả bờ là đường thẳng $2  x-3 y+5=0$ chứa điểm $O$.
	 }
\end{ex}

% nguồn: MaHoaNH_TracNghiem_Lop10
\begin{ex}%[0D2H1-2]%[ma:5f1fd454]
	Cho bất phương trình bậc nhất hai ẩn $x-2 y+2 \le 0$. Chọn khẳng định đúng trong các khẳng định dưới đây? 
	\choice 
		{\True $(0;0)$ không là một nghiệm của bất  phương trình}
		{$(1;4)$ không là  một nghiệm của bất phương trình}
		{$(0; 3)$ không là một nghiệm của bất phương trình} 
		{$(2; 2)$ không  là một nghiệm cùa bất phương trình}
	 \loigiai{
		\begin{itemize}
			\item Xét  cặp số $(0; 0)$, ta có $0 -2\cdot 0 +2 = 2 >0\Rightarrow (0; 0)$ không là một nghiệm của bất phương trình $x-2 y+2 \le 0$.
			\item Xét cặp số $(1;4)$, ta có $1 -2\cdot 4 +2 = -5 <0\Rightarrow (1;4)$ là một nghiệm của bất phương trình  $x-2 y+2 \le 0$.
			\item Xét cặp số $(0; 3)$, ta có $0 -2\cdot 3 +2 = -4 <0\Rightarrow (0; 3)$ là một nghiệm của bất phương trình $x-2  y+2 \le 0$.
			\item Xét cặp số $(2; 2)$, ta  có $2 -2\cdot 2 +2 = 0\Rightarrow (2; 2)$ là một nghiệm của bất phương trình $x-2 y+2 \le 0$.
		\end{itemize}
		Vậy khẳng định $(0;0)$ không là nghiệm của bất phương trình $x-2 y+2 \le 0$ là khẳng định đúng. 
	} 
\end{ex}

% nguồn: MaHoaNH_TracNghiem_Lop10
\begin{ex}%[0D2H1-2]%[ma:818056a8]
	Cho bất phương  trình $-2 x+3 y>3$. Khẳng định nào dưới đây \textbf{sai}?
	\choice 
		{$(0; 0)$ không là  một nghiệm bất phương trình}
		{\True $(-1; 1)$ không  là một nghiệm bất phương trình}
		{$(0; 1)$  không là một nghiệm bất phương trình}
		{$(1; 3)$ là một  nghiệm bất phương trình}
	 \loigiai{
		\begin{itemize}
			\item Xét cặp số $(0; 0)$, ta có $-2\cdot0 +3\cdot 0 = 0 <3\Rightarrow (0; 0)$ không là một nghiệm của bất phương trình $-2 x+3  y>3$.
			\item Xét cặp số $(-1; 1)$, ta có $-2\cdot(-1) +3\cdot 1 = 5 >3\Rightarrow (-1; 1)$ là  một nghiệm của bất phương trình $-2 x+3 y>3$.
			\item Xét cặp số $(0; 1)$, ta có $-2\cdot0 +3\cdot 1 = 3 \Rightarrow (0; 1)$ không là một nghiệm của bất phương trình  $-2 x+3 y>3$.
			\item Xét cặp số $(1; 3)$, ta  có $-2\cdot1 +3\cdot 3 = 7 >3\Rightarrow (1; 3)$ là một nghiệm của bất phương trình $-2 x+3 y>3$.
		\end{itemize}
		Vậy khẳng định $(-1; 1)$ không là một nghiệm bất phương trình  $-2 x+3 y>3$ là khẳng định sai.
	 }
\end{ex}

% nguồn: MaHoaNH_TracNghiem_Lop10
\begin{ex}%[0D2H1-2]%[ma:bf3b8fbf]
	Điểm $O(0; 0)$ thuộc miền nghiệm  của bất phương trình nào dưới đây?
	 \choice
		{$x+3 y+2 \le 0$} 
		{$x+y+2  \le 0$}
		 {\True $2 x+5 y-2<0$}
		{$2 x+y<0$} 
	 \loigiai{
		\begin{itemize}
			\item Xét bất phương trình $x+3 y+2 \le  0$, ta có $0 + 3\cdot 0 + 2 = 2>0$. Do đó điểm $O(0; 0)$ không thuộc miền nghiệm của bất phương trình $x+3 y+2 \le 0$.
			\item  Xét bất phương trình $x+y+2 \le 0$, ta có $0 + 0 + 2 = 2>0$. Do đó điểm $O(0; 0)$ không thuộc miền nghiệm của bất phương trình $x+3 y+2 \le 0$.
			\item Xét bất phương trình $2 x+5 y-2<0$, ta có $2\cdot0 + 5\cdot 0 -2=-2  < 0$. Do đó điểm $O(0; 0)$ thuộc miền nghiệm của bất phương trình $2 x+5 y-2<0$.
			\item Xét bất phương trình $2 x+y<0$, ta có $2\cdot0 + 0 = 0$. Do đó điểm $O(0; 0)$ không thuộc  miền nghiệm của bất phương trình $2 x+y<0$.
		\end{itemize}
		Vậy điểm  $O(0; 0)$ thuộc miền nghiệm của bất phương trình $2 x+5 y-2<0$.
	 }
\end{ex}

% nguồn: MaHoaNH_TracNghiem_Lop10
\begin{ex}%[0D2H1-2]%[ma:e2295743]
	Cho bất phương trình bậc nhất hai ẩn $2 x-5 y \leq 8$. Khẳng định nào dưới đây là khẳng  định \textbf{sai}?
	\choice 
		{$(3;-4)$ không là một nghiệm của bất phương  trình}
		 {\True $(-2; 2)$ không là một nghiệm của bất phương trình}
		{$(-3;-1)$  là một nghiệm của bất phương trình}
		{$(5; 0)$ không là một nghiệm  của bất phương trình}
	 \loigiai{
		\begin{itemize}
			\item Xét cặp số $(3;-4)$, ta có $2\cdot3  -5\cdot (-4) = 26 >8\Rightarrow (3;-4)$ không là một nghiệm của bất phương trình $2 x-5 y \leq 8$.
			\item Xét cặp số $(-2; 2)$, ta có $2\cdot(-2) -5\cdot 2 = -14 \le 8\Rightarrow (-2;2)$ là một nghiệm  của bất phương trình $2 x-5 y \leq 8$.
			\item Xét cặp số $(-3;-1)$, ta có $2\cdot(-3) -5\cdot (-1) =-1 \le 8\Rightarrow (-3;-1)$ là một nghiệm của bất phương trình $2  x-5 y \leq 8$.
			\item  Xét cặp số $(5; 0)$, ta có $2\cdot5 -5\cdot 0 = 10 >8\Rightarrow (5; 0)$ không là một nghiệm của bất phương trình $2 x-5 y \leq 8$.
		\end{itemize}
		Vậy khẳng định $(-2; 2)$ không là một nghiệm của bất phương trình  $2 x-5 y \leq 8$ là khẳng định sai.
	} 
\end{ex}

% nguồn: MaHoaNH_TracNghiem_Lop10
\begin{ex}%[0D2H1-2]%[ma:41db1d43]
	Cho bất phương trình $4 x-3  y \le 5$. Khẳng định nào dưới đây đúng?
	\choice 
		{$(1;-1)$ là một nghiệm của bất phương trình} 
		{$(0; 0)$ không là một nghiệm của bất phương  trình}
		{\True $(2; 1)$ là một nghiệm của bất phương  trình}
		{$(3;-1)$ là một  nghiệm của bất phương trình}
	 \loigiai{
		\begin{itemize}
			 \item Xét cặp số $(1;-1)$, ta có $4\cdot 1 -3\cdot(-1) = 7 >5\Rightarrow (1;-1)$ không là một nghiệm của bất phương trình $4 x-3 y \le 5$.
			\item Xét cặp số $(0; 0)$, ta có $4\cdot 0 -3\cdot0 = 0 <5\Rightarrow  (0; 0)$ là một nghiệm của bất phương trình $4 x-3 y \le 5$.
			\item Xét cặp số  $(2; 1)$, ta có $4\cdot2 -3\cdot1=5\Rightarrow (2; 1)$ là một nghiệm của bất phương trình $4 x-3 y \le 5$.
			\item Xét cặp số $(3;-1)$, ta có $4\cdot3-3\cdot(-1)=15>5\Rightarrow (3;-1)$ không là một nghiệm của  bất phương trình $4 x-3 y \le 5$.
		\end{itemize}
		Vậy khẳng định $(2; 1)$ là một nghiệm của bất phương trình $4  x-3 y \le 5$ là khẳng định đúng.
	} 
\end{ex}

% nguồn: MaHoaNH_TracNghiem_Lop10
\begin{ex}%[0D2H1-2]%[ma:66938b6d]
	Có bao nhiêu  giá trị nguyên dương của tham số $m$ để nghiệm của hệ $
	\heva{&2x + y = 3 \\&x + y = 4}$ là nghiệm của bất phương trình $(m-1)x + 3y  + 1 \ge 0$.
	 \choice
		{$15$} 
		{$16$} 
		 {$18$}
		 {\True $17$}
	 \loigiai{
		Ta có $\heva{&2x + y = 3 \\&x + y =  4}
		 \Leftrightarrow
		\heva{&x  = -1 \\&y = 5}\Rightarrow A(-1;5)$.\\
		Yêu cầu của bài toán  tương đương điểm $A(-1;5)$ thuộc miền nghiệm của bất phương trình $(m-1)x + 3y + 1 \ge 0$, tức là ta có
		$$-m + 1 + 15  + 1 \ge 0\Leftrightarrow m \le 17.$$
		Vì $m$ có giá trị nguyên dương nên có $17$ giá trị của  $m$ thỏa mãn.
	 }
\end{ex}

% nguồn: MaHoaNH_TracNghiem_Lop10
\begin{ex}%[0D2H1-2]%[ma:6ae13cd4]
	 \immini[thm]{
	Nửa mặt phẳng kể cả bờ (phần không bị gạch) trong hình vẽ dưới đây là miền nghiệm của bất phương  trình nào?
	 \choice
		{\True $2x+y\ge  2$}
		 {$2x+y\le 2$}
		 {$2x-y\ge 2$}
		 {$2x-y\le 2$}}
		{	\begin{tikzpicture}[>=stealth,x=1cm,y=1cm,scale=.7]
		\draw[->] (0,-4)--(0,4)  node[left]{$y$};
		\draw (0,0) node[below left]{$O$}; 
		 \path[pattern=horizontal lines] (-1,4)-|(-4,-4)--(3,-4)--cycle;
		\draw  (-1,4)--(3,-4);
		\fill  (0,2) circle (1pt) node[right]{$2$}
			(1,0) circle (1pt)  node[above]{$1$}
			(0,0) circle  (1pt);
		 \draw[->] (-4,0)--(4,0) node[below right]{$x$}; %[0D2N2-1]
	\end{tikzpicture}
		} 
	\loigiai{ 
		\begin{itemize}
			 \item Cách 1:\\
			Ta thấy bờ của nửa mặt phẳng đã cho là đường thẳng $d$ cắt các trục tọa độ tại $A\left(0;2\right), B\left(1;0\right)$  nên $d$ có phương trình $y=-2x+2$ $\Leftrightarrow 2x+y=2$.\\
			 Vì điểm $O\left(0;0\right)$ không thuộc miền nghiệm của bất phương trình nên bất phương trình cần tìm là $2x+y\ge 2$.
			 \item Cách 2:\\
			Xét điểm $O$ không thuộc miền nghiệm $\Rightarrow$ loại $2x+y\le 2$,  $2x-y\le 2$.\\
			Xét  điểm $A\left(0;2\right)$ thuộc miền nghiệm $\Rightarrow$ loại $2x-y\ge 2$.
	\end{itemize}}
\end{ex}

% nguồn: MaHoaNH_TracNghiem_Lop10
\begin{ex}%[0D2H1-2]%[ma:68e1172e]
	Đường thẳng $d\colon 2x-y=2$ chia mặt phẳng tọa độ thành hai miền I, II có bờ là đường thẳng $d$. Xác định miền nghiệm của bất phương  trình $2x-y\ge 2$.
	\begin{center}
		\begin{tikzpicture}[line join=round,line cap=round,>=stealth]
				\tikzset{every  node/.style={scale=0.9}}
				\draw[->] (-1.1,0)--(3.1,0) node[below left]  {$x$};
				 \draw[->] (0,-2.5)--(0,3.1) node[below left] {$y$};
				\draw (0,0)  node [below left] {$O$};
				 \foreach \x/\nx in {1/1,2/2}
				\draw[thin]  (\x,1pt)--(\x,-1pt) node [below] {$\nx$};
				 \foreach \y/\ny in {-2/-2,-1/-1,1/1,2/2}
				\draw[thin] (1pt,\y)--(-1pt,\y) node [left] {$\ny$}; 
				\begin{scope}
					\clip (-1,-2.5)  rectangle (3,3);
					 \draw[samples=200,domain=-1:3,smooth,variable=\x] plot (\x,{2*(\x)+-2});
				\end{scope}
				 \node [right,rotate=65] at (1,.5){$2x-y=2$};
				\path (0.5,1.5) node [above]{II}; 
				\path (2,0.5)  node [above]{I};
			\end{tikzpicture}
	\end{center}
	 \choice
		{Nửa  mặt phẳng I không kể đường thẳng $d$}
		{\True Nửa mặt phẳng I kể cả đường  thẳng $d$}
		{Nửa  mặt phẳng II kể cả đường thẳng $d$}
		{Nửa mặt phẳng II  không kể đường thẳng $d$}
	 \loigiai{
		Thay tọa độ điểm $O(0;0)$ vào bất phương trình đã cho $2x-y\ge 2$, ta có $2\cdot0 -0\ge 2$ là mệnh đề sai. Do vậy  miền nghiệm của bất phương trình đã cho là miền I không chứa điểm $O$, kể cả bờ $d$.
	 }
\end{ex}

% nguồn: MaHoaNH_TracNghiem_Lop10
\begin{ex}%[0D2H1-2]%[ma:257dd032]
	Hình dưới đây biểu diễn  hình học miền nghiệm của bất phương trình nào? (Miền nghiệm là miền không gạch chéo và miền nghiệm không chứa đường thẳng)
	\begin{center}
		\begin{tikzpicture}[line join=round,line cap=round,>=stealth,thick]
				 \tikzset{every node/.style={scale=0.9}}
				\begin{scope}
					\clip (-2,-2) rectangle (4,4); 
					 \fill[pattern=north east lines,,opacity=0.5] (-2,-2)--(2,-2)--(2,4)--cycle;
					\draw (-2,-2)--(2,4) node  [pos=0.9,left] {$(d)$};
				\end{scope}
				 \draw[->] (-2,0)--(2,0) node[below]{$x$};
				\draw[->]  (0,-2)--(0,4) node[left]{$y$};
				\draw (0,0) node[below  left]{$O$};
				\foreach \x  in {-1,1}
				\draw[thin]  (\x,1pt)--(\x,-1pt) node [above] {$\x$};
				 \foreach \y in {1}
				\draw[thin] (1pt,\y)--(-1pt,\y)  node [left] {$\y$};
				\draw[dashed,thin] (0,0)--(0,2)--(0,2); 
				 \draw[dashed,thin] (1,0)--(1,0)--(0,0);
			\end{tikzpicture}
	\end{center}
	 \choice
		 {$3x+2y>2$}
		 {$3x+2y<2$}
		 {\True $-3x+2y>2$}
		 {$-3x+2y<2$}
	 \loigiai{
		 Đường thẳng trong hình vẽ là $-3x+2y=2$.\\
		Gốc tọa độ $O(0;0)$ không thuộc miền  nghiệm nên ta chọn đáp án$-3x+2y>2$.
	} 
\end{ex}

% nguồn: MaHoaNH_TracNghiem_Lop10
\begin{ex}%[0D2H1-2]%[ma:164defe0]
	\immini{Nửa mặt phẳng không bị gạch (không kể  bờ $d$) ở hình bên là miền nghiệm của bất phương trình nào đây?
	 \choice
		 {$3x+y<3$}
		{\True  $3x+y>3$}
		 {$x+3y<3$}
		{$x+3y>3$}} 
		{	\begin{tikzpicture}[scale=1,line join=round,line cap=round,>=stealth,thick]
		\tikzset{every node/.style={scale=0.9}} 
		\begin{scope}
			\clip (-2,-1) rectangle  (3,4);
			\fill[pattern=north east lines] (-2,9)--(-2,-6)--(3,-6)--cycle; 
			\draw  (-0.33,4)--(1.33,-1) node [pos=0.45,above,sloped] {$d$};
		\end{scope}
		\draw[->] (-2,0)--(3,0) node[below]{$x$}; 
		\draw[->] (0,-1)--(0,4)  node[left]{$y$};
		\draw (0,0) node[below  left]{$O$};
		\foreach \x  in {1}
		\draw[thin] (\x,1pt)--(\x,-1pt)  node [above] {$\x$};
		\foreach \y  in {3}
		\draw[thin] (1pt,\y)--(-1pt,\y) node  [right] {$\y$};
	\end{tikzpicture}
		 }
	\loigiai{ 
		Gọi đường thẳng $d\colon y=ax+b$ đi qua hai điểm $(1;0)$ và $(0;3)$.\\ 
		Suy ra $\heva{&a+b=0\\&b=3}\Rightarrow\heva{&a=-3\\&b=3}\Rightarrow d\colon y=-3x+3\Leftrightarrow 3x+y=3$.  \\
		 Xét điểm $O(0;0)\notin d$, ta có $3\cdot0 +0=0<3$.\\
		Ta thấy điểm $(0;0)$  không thuộc miền nghiệm (phần không bị gạch).\\
		Do đó $O(0;0)$ không thuộc miền nghiệm của bất  phương trình $3x+y>3$.\\
		Vậy nửa mặt phẳng phần không gạch sọc (không  chứa điểm $O(0; 0)$) là miền nghiệm của bất phương trình $3x+y>3$.
	} 
\end{ex}

% nguồn: MaHoaNH_TracNghiem_Lop10
\begin{ex}%[0D2H1-2]%[ma:2bf3c9e1]
	\immini{Nửa mặt phẳng phần  không gạch sọc được cho bởi hình bên là miền nghiệm của bất phương trình nào dưới đây?
	\choice 
		 {$2x+y-6>0$}
		{\True  $2x+y-6<0$}
		{$x+2y-6<0$} 
		{$x+2y-6>0$} 
		} 
		{ 
		\begin{tikzpicture}[line join=round,line cap=round,>=stealth,scale=.7]
				\tikzset{every  node/.style={scale=0.9}}
				\begin{scope}
					\clip (-2,-1) rectangle (5,7); 
					\fill[pattern=north east lines]  (-3,12)--(6,12)--(6,-6)--cycle;
					\draw[samples=200,smooth,variable=\x]  plot (\x,{-2*(\x)+6});
				\end{scope}
				\draw[->]  (-2,0)--(5.2,0) node[below]{$x$};
				\draw[->] (0,-1)--(0,7.2) node[left]{$y$}; 
				\draw  (0,0) node[below left]{$O$};
				\foreach \x in {3} 
				\draw[thin] (\x,1pt)--(\x,-1pt) node [below]  {$\x$};
				\foreach \y in {6} 
				\draw[thin] (1pt,\y)--(-1pt,\y) node  [left] {$\y$};
			\end{tikzpicture}
		 }
	 \loigiai{
		Gọi đường thẳng $d\colon y=ax+b$ đi qua hai  điểm $(3;0)$ và $(0;6)$.\\
		Suy ra $\heva{&3a+b=0\\&b=6}\Rightarrow\heva{&a=-2\\&b=6}\Rightarrow d\colon y=-2x+6\Leftrightarrow  2x+y-6=0.$\\
		Xét điểm $O(0;0)\notin d$,  ta có $2\cdot0 +0-6=-6<0$.\\
		Ta thấy điểm $O(0; 0)$ thuộc miền  nghiệm (phần phông gạch sọc).\\
		Do đó $O(0;0)$ thuộc miền nghiệm của bất phương  trình $2x+y-6<0$.\\
		Vậy nửa mặt phẳng phần không gạch sọc (chứa điểm $O$) là miền nghiệm của bất phương trình  $2x+y-6<0$.
	 }
\end{ex}

% nguồn: MaHoaNH_TracNghiem_Lop10
\begin{ex}%[0D2H1-2]%[ma:76791e6a]
	Trong hệ tọa độ $Oxy$, cho bất phương trình $2x+y \geq 2$ có miền nghiệm $D$. Dựng hình vuông $ABCO$ có cạnh $a$ nằm trong góc phần tư thứ nhất, với $O(0; 0)$ là gốc tọa độ. Biết rằng diện tích phần chung giữa miền nghiệm $D$ và hình vuông $ABCO$ bằng $2\,022$. Khi đó giá trị  của $a$ thuộc khoảng nào trong các khoảng sau đây?
	\choice 
		{\True  $(44; 45)$}
		{$(44;  44,5)$}
		{$(45; 46)$} 
		{$(43; 44)$} 
	\loigiai{ 
	\immini{Vẽ đường thẳng $d\colon 2x+y=2$.\\ 
		Thay điểm $O(0; 0)$ vào bất phương trình $2x+y \geq 2$ ta được $0\geq 2$ (vô lí) nên điểm $O(0; 0)$ không thuộc miền nghiệm $D$ của bất phương trình $2x+y \geq  2$.\\
		Khi đó miền nghiệm $D$ của bất phương trình $2x+y \geq 2$ là miền không chứa điểm $O(0; 0)$ có bờ $d$, kể cả  đường thẳng $d$ (miền không gạch như hình vẽ).\\
		Vẽ hình vuông $ABCO$ cạnh bằng $a$ nằm trong góc phần tư thứ nhất trên cùng hệ trục tọa độ với miền  nghiệm $D$.}
		{{\begin{tikzpicture}[scale=0.9,font=\footnotesize,line join=round,line cap=round,>=stealth]
		\draw[->]  (-1.5,0)--(5,0) node[below] {$x$};
		\draw[->] (0,-2)--(0,5) node[left]  {$y$};
		 \fill (0,0) circle (1pt) node[below left]{$O$};
		\foreach \x in {-1,1,2,3,4} \draw (\x,0.1)--(\x,-0.1); 
		\foreach \y in {-1,1,2,3,4}  \draw (0.1,\y)--(-0.1,\y);
		 \path
			(1,0)  coordinate (F)
			(4,0) coordinate  (A)
			(0,4) coordinate (C) 
			(4,4) coordinate (B) 
			(0,2) coordinate  (E);
		\foreach \x/\g in {A/-90,B/45,C/180,E/0,F/-45} 
		\fill (\x)  circle (1pt)
			 ($(\g:3mm)+(\x)$) node {$\x$};
		\draw (E)--(F) 
			(C)--(B)--(A); 
		 \fill;
		\fill[pattern=north east lines,smooth] 
			(-1.5,5)--(2,-2)--(-1.5,-2)--(-1.5,5); 
		\draw[thick,smooth,samples=200] plot[domain=-1.5:2] (\x,{-2*(\x)+2}); 
	\end{tikzpicture}
		 }}\noindent
		Dựa vào hình vẽ ta thấy diện tích phần chung giữa miền  nghiệm $D$ và hình vuông (phần tô màu) là $S=S_{ABCO}-S_{OEF}=OA^2-\dfrac{1}{2}\cdot OE\cdot OF=a^2-\dfrac{1}{2}\cdot 1\cdot 2=a^2-1$.\\
		Mặt khác $S=2\,022$ nên $a^2-1=2\,022\Leftrightarrow a^2=2\,023\Rightarrow a=\sqrt{2\,023} \Rightarrow a \in(44; 45)$. 
	 }
\end{ex}

% nguồn: MaHoaNH_TracNghiem_Lop10
\begin{ex}%[0D2H1-2]%[ma:c5b0f9cb]
	Miền nghiệm của bất phương trình $3x-2y >-6$ là miền không bị gạch (không tính đường  biên) trong hình nào sau đây?
	 \choice
		{\begin{tikzpicture}[>=stealth,x=1cm,y=1cm,scale=.6]
		 \draw[->] (-2,0)--(4,0) node[below right]{$x$};
		\draw[->]  (0,-3)--(0,4) node[left]{$y$};
		 \draw (0,0) node[below left]{$O$};
		 \path[pattern=north west lines] (-0.67,4)-|(-2,-3)-|(4,-3)--cycle;
		 \draw (-0.67,4)--(4,-3);
		\fill (2,0) circle (1pt) node[above  right]{$2$}
			(0,3)  circle (1pt) node[right]{$3$};
	\end{tikzpicture}
		} 
		{\begin{tikzpicture}[>=stealth,x=1cm,y=1cm,scale=.6]
		\draw[->]  (-4,0)--(2,0) node[below right]{$x$};
		\draw[->] (0,-3)--(0,4) node[left]{$y$}; 
		\draw (0,0) node[below  left]{$O$};
		\path[pattern=north east lines]  (-4,-3)|-(2,-3)|-(0.67,4)--cycle;
		\draw (-4,-3)--(0.67,4); 
		\fill (-2,0)  circle (1pt) node[above left]{$-2$}
			(0,3) circle  (1pt) node[left]{$3$};
	\end{tikzpicture}
		}{\True \begin{tikzpicture}[>=stealth,x=1cm,y=1cm,scale=.6]
		\draw[->] (0,-3)--(0,4) node[left]{$y$}; 
		\draw (0,0)  node[below left]{$O$};
		\path[pattern=horizontal lines]  (-4,-3)|-(0.67,4)--cycle;
		\draw  (-4,-3)--(0.67,4);
		 \draw[->] (-4,0)--(2,0) node[below right]{$x$};
		\fill (-2,0)  circle (1pt) node[below right]{$-2$}
			(0,3)  circle (1pt) node[right]{$3$};
	\end{tikzpicture}
		} 
		{\begin{tikzpicture}[>=stealth,x=1cm,y=1cm,scale=.6]
		\draw[->] (-4,0)--(2,0)  node[below right]{$x$};
		\draw[->]  (0,-3)--(0,4) node[left]{$y$};
		 \draw (0,0) node[below left]{$O$};
		\path[pattern=vertical lines]  (-4,3)|-(0.6,-4)--cycle;
		\draw (-4,3)--(0.67,-4); 
		\fill  (-2,0) circle (1pt) node[above]{$-2$}
			(0,-3) circle  (1pt) node[right]{$-3$};
	\end{tikzpicture}
	} 
	\loigiai{ 
		\begin{center}
			\begin{tikzpicture}[>=stealth,x=1cm,y=1cm,scale=.6]
					\draw[->]  (0,-3)--(0,4) node[left]{$y$};
					 \draw (0,0) node[below left]{$O$};
					 \path[pattern=horizontal lines] (-4,-3)|-(0.67,4)--cycle;
					 \draw (-4,-3)--(0.67,4);
					\draw[->] (-4,0)--(2,0) node[below  right]{$x$};
					\fill  (-2,0) circle (1pt) node[below right]{$-2$}
						(0,3) circle (1pt)  node[right]{$3$};
				\end{tikzpicture}
		\end{center}
		Trước hết, ta vẽ đường thẳng  $(d)\colon 3x-2y=-6$.\\
		Ta thấy $\left(0; 0\right)$ là một nghiệm  của bất phương trình đã cho.\\
		Vậy miền nghiệm  cần tìm là nửa mặt phẳng bờ $(d)$ chứa điểm $\left(0; 0\right)$.
	 }
\end{ex}

% nguồn: MaHoaNH_TracNghiem_Lop10
\begin{ex}%[0D2H1-2]%[ma:edc5ab6b]
	Miền nghiệm (miền không gạch sọc) của  bất phương trình $3x+y\ge 3$ là
	 \choice
		{\True \begin{tikzpicture}[line join=round,line cap=round,>=stealth,thick]
		\tikzset{every  node/.style={scale=0.9}}
		\begin{scope}
			\clip (-1,-1)  rectangle (2,4);
			\fill[pattern=north  east lines]
				 (-1/3,4)--(-1,4)--(-1,-1)--(4/3,-1)--cycle;
			\draw (-1,6)--(2,-3); 
		\end{scope}
		\draw[->] (-1,0)--(2,0) node[below]{$x$}; 
		\draw[->]  (0,-1)--(0,4) node[left]{$y$};
		\draw  (0,0) node[below left]{$O$};
		\fill[black] 
			 (0,3) node[right]{$3$}circle (1pt)
			(1,0)  node[above right]{$1$}circle (1pt);
	\end{tikzpicture}
		} 
		{\begin{tikzpicture}[line join=round,line cap=round,>=stealth,thick]
		\tikzset{every  node/.style={scale=0.9}}
		\begin{scope}
			\clip  (-2,-1) rectangle (1,4);
			\fill[pattern=north east  lines]
				 (1/3,4)--(1,4)--(1,-1)--(-4/3,-1)--cycle;
			\draw  (1,6)--(-2,-3);
		\end{scope}
		\draw[->]  (-2,0)--(1,0) node[below]{$x$};
		 \draw[->] (0,-1)--(0,4) node[left]{$y$};
		\draw (0,0)  node[above right]{$O$};
		\fill[black] 
			(0,3) node[left]{$3$}circle  (1pt)
			(-1,0) node[above  left]{$-1$}circle (1pt);
	\end{tikzpicture}
		 }
		{\begin{tikzpicture}[line join=round,line cap=round,>=stealth,thick]
		\tikzset{every  node/.style={scale=0.9}}
		\begin{scope}
			\clip (-2,-1)  rectangle (1,4);
			\fill[pattern=north  east lines]
				(1/3,4)--(-2,4)--(-2,-1)--(-4/3,-1)--cycle; 
			\draw  (1,6)--(-2,-3);
		\end{scope}
		\draw[->]  (-2,0)--(1,0) node[below]{$x$};
		\draw[->]  (0,-1)--(0,4) node[left]{$y$};
		\draw (0,0) node[above right]{$O$}; 
		 \fill[black]
			(0,3) node[right]{$3$}circle  (1pt)
			(-1,0)  node[below right]{$-1$}circle (1pt);
	\end{tikzpicture}
		 }
		{\begin{tikzpicture}[line join=round,line cap=round,>=stealth,thick]
		\tikzset{every  node/.style={scale=0.9}}
		\begin{scope}
			 \clip (-2,1) rectangle (1,-4);
			 \fill[pattern=north east lines]
				(-4/3,1)--(-2,1)--(-2,-4)--(1/3,-4)--cycle; 
			 \draw (1,-6)--(-2,3);
		\end{scope}
		\draw[->]  (-2,0)--(1,0) node[below]{$x$};
		\draw[->]  (0,-4)--(0,1) node[left]{$y$};
		\draw (0,0)  node[above right]{$O$};
		 \fill[black]
			(0,-3) node[right]{$-3$}circle  (1pt)
			(-1,0) node[above  right]{$-1$}circle (1pt);
	\end{tikzpicture}
		 }
	 \loigiai{
	Đường thẳng đi qua điểm có tọa độ $(0;3)$, $(1;0)$, mặt khác tọa độ điểm $O$ không thuộc miền nghiệm nên chọn phương án: Nửa mặt phẳng bờ là đường  thẳng $d\colon3x+y=3$ (kể cả bờ $d$) chứa điểm $M(2;1)$.}
\end{ex}

% nguồn: MaHoaNH_TracNghiem_Lop10
\begin{ex}%[0D2H1-2]%[ma:6fdbff62]
	Miền nghiệm của bất phương trình $-x+y < 2$ được xác  định bởi miền nào (nửa mặt phẳng không bị gạch và không kể đường thẳng $d$) sau đây?
	 \choice
		{\begin{tikzpicture}[line join=round,line cap=round,>=stealth,thick]
		\tikzset{every node/.style={scale=0.9}} 
		\begin{scope}
			\clip (-2,-1) rectangle (3,4); 
			\fill[pattern=north  east lines]
				 (-2,4)--(-2,-1)--(3,-1)--cycle;
			\draw  (-2,4)--(3,-1);
		\end{scope}
		\draw[->] (-2,0)--(3,0)  node[below]{$x$};
		 \draw[->] (0,-1)--(0,4) node[left]{$y$};
		\draw (0,0) node[below  left]{$O$};
		 \fill[black]
			(0,2)  node[right]{$2$}circle (1pt)
			(2,0) node[above]{$2$}circle (1pt); 
	\end{tikzpicture}
		 }
		{\begin{tikzpicture}[line join=round,line cap=round,>=stealth,thick]
		\tikzset{every  node/.style={scale=0.9}}
		\begin{scope}
			\clip (-2,-1) rectangle (3,4); 
			 \fill[pattern=north east lines]
				(-2,4)--(3,4)--(3,-1)--cycle; 
			 \draw (-2,4)--(3,-1);
		\end{scope}
		 \draw[->] (-2,0)--(3,0) node[below]{$x$};
		 \draw[->] (0,-1)--(0,4) node[left]{$y$};
		\draw  (0,0) node[below left]{$O$};
		 \fill[black]
			(0,2)  node[below left]{$2$}circle (1pt)
			(2,0) node[below]{$2$}circle  (1pt);
	\end{tikzpicture}
		} 
		{\begin{tikzpicture}[line join=round,line cap=round,>=stealth,thick]
		 \tikzset{every node/.style={scale=0.9}}
		\begin{scope}
			\clip  (-3,-1) rectangle (1,3);
			\fill[pattern=north west  lines]
				(-3,-1)--(1,-1)--(1,3)--cycle; 
			\draw (-3,-1)--(2,4); 
		\end{scope}
		 \draw[->] (-3,0)--(1,0) node[below]{$x$};
		\draw[->] (0,-1)--(0,3)  node[left]{$y$};
		 \draw (0,0) node[above right]{$O$};
		\fill[black] 
			(0,2) node[left]{$2$}circle (1pt) 
			(-2,0)  node[above left]{$-2$}circle (1pt);
	\end{tikzpicture}
		 }
		{\True \begin{tikzpicture}[line join=round,line cap=round,>=stealth,thick]
		\tikzset{every  node/.style={scale=0.9}}
		\begin{scope}
			 \clip (-3,-1) rectangle (1,3);
			\fill[pattern=north  west lines]
				(-3,-1)--(-3,3)--(1,3)--cycle; 
			\draw (-3,-1)--(2,4); 
		\end{scope}
		\draw[->]  (-3,0)--(1,0) node[below]{$x$};
		\draw[->] (0,-1)--(0,3)  node[left]{$y$};
		\draw (0,0) node[above right]{$O$}; 
		 \fill[black]
			(0,2) node[below right]{$2$}circle  (1pt)
			 (-2,0) node[below]{$-2$}circle (1pt);
	\end{tikzpicture}
		} 
	 \loigiai{
	\immini{Ta có đường thẳng $\Delta\colon -x+y =2$ đi qua hai điểm $(0;2)$ và $(-2;0)$ và điểm $O(0;0)$ thoả mãn bất phương trình $-x+y < 2$ nên miền nghiệm là nửa mặt phẳng bờ $\Delta$ (không kể bờ) chưa điểm $O$ là miền nghiệm. Đáp án là hình vẽ bên}{\begin{tikzpicture}[line join=round,line cap=round,>=stealth,thick]
		\tikzset{every node/.style={scale=0.9}} 
		\begin{scope}
			 \clip (-3,-1) rectangle (1,3);
			\fill[pattern=north  west lines]
				(-3,-1)--(-3,3)--(1,3)--cycle; 
			 \draw (-3,-1)--(2,4);
		\end{scope}
		\draw[->]  (-3,0)--(1,0) node[below]{$x$};
		\draw[->]  (0,-1)--(0,3) node[left]{$y$};
		\draw  (0,0) node[above right]{$O$};
		 \fill[black]
			 (0,2) node[below right]{$2$}circle (1pt)
			(-2,0) node[below]{$-2$}circle  (1pt);
	\end{tikzpicture}
		 }
	 }
\end{ex}

% nguồn: MaHoaNH_TracNghiem_Lop10
\begin{ex}%[0D2H1-2]%[ma:a5a2fa85]
	Biểu diễn miền nghiệm của bất phương trình $x-y<3$ trên mặt phẳng tọa độ $Oxy$ là phần không  gạch sọc trong hình nào sau đây?
	 \choice
		{\begin{tikzpicture}[line join=round,line cap=round,>=stealth,thick,scale=0.8]
		 \tikzset{every node/.style={scale=0.9}}
		\begin{scope}
			\clip  (-1,-1) rectangle (4,4);
			 \fill[pattern=north east lines] (-2,5)--(-2,-2)--(5,-2)--cycle;
			 \draw (-1,4)--(4,-1);
		\end{scope}
		\draw[->] (-1,0)--(4,0)  node[below]{$x$};
		\draw[->]  (0,-1)--(0,4) node[left]{$y$};
		\draw  (0,0) node[below left]{$O$};
		 \foreach \x in {3}
		 \draw[thin] (\x,1pt)--(\x,-1pt) node [above] {$\x$};
		\foreach \y in  {3}
		\draw[thin] (1pt,\y)--(-1pt,\y) node  [right] {$\y$};
	\end{tikzpicture}
		 }
		{ 
		\begin{tikzpicture}[line join=round,line cap=round,>=stealth,thick,scale=0.8]
				\tikzset{every node/.style={scale=0.9}} 
				\begin{scope}
					\clip  (-1,-1) rectangle (4,4);
					\fill[pattern=north east lines]  (-2,5)--(5,5)--(5,-2)--cycle;
					\draw  (-1,4)--(4,-1);
				\end{scope}
				\draw[->] (-1,0)--(4,0) node[below]{$x$}; 
				 \draw[->] (0,-1)--(0,4) node[left]{$y$};
				 \draw (0,0) node[below left]{$O$};
				\foreach \x  in {3}
				 \draw[thin] (\x,1pt)--(\x,-1pt) node [below] {$\x$};
				\foreach  \y in {3}
				\draw[thin] (1pt,\y)--(-1pt,\y) node [left] {$\y$}; 
			\end{tikzpicture}
		} 
		{ 
		\begin{tikzpicture}[line join=round,line cap=round,>=stealth,thick,scale=0.8]
				\tikzset{every  node/.style={scale=0.9}}
				\begin{scope}
					 \clip (-1,-4) rectangle (4,1);
					\fill[pattern=north west lines] (-2,-5)--(-2,2)--(5,2)--cycle; 
					\draw (4,1)--(-1,-4); 
				\end{scope}
				 \draw[->] (-1,0)--(4,0) node[below]{$x$};
				\draw[->] (0,-4)--(0,1)  node[left]{$y$};
				 \draw (0,0) node[below left]{$O$};
				 \foreach \x in {3}
				\draw[thin] (\x,1pt)--(\x,-1pt) node [below] {$\x$}; 
				\foreach \y in {-3} 
				\draw[thin] (1pt,\y)--(-1pt,\y) node [right] {$\y$}; 
			\end{tikzpicture}
		} 
		 {\True
		\begin{tikzpicture}[line join=round,line cap=round,>=stealth,thick,scale=0.8]
				\tikzset{every  node/.style={scale=0.9}}
				\begin{scope}
					\clip (-1,-4) rectangle (4,1); 
					 \fill[pattern=north west lines] (-2,-5)--(5,-5)--(5,2)--cycle;
					\draw  (4,1)--(-1,-4);
				\end{scope}
				 \draw[->] (-1,0)--(4,0) node[below]{$x$};
				 \draw[->] (0,-4)--(0,1) node[left]{$y$};
				 \draw (0,0) node[below left]{$O$};
				\foreach \x  in {3}
				\draw[thin]  (\x,1pt)--(\x,-1pt) node [above] {$\x$};
				 \foreach \y in {-3}
				\draw[thin] (1pt,\y)--(-1pt,\y) node  [left] {$\y$};
			\end{tikzpicture}
		 }
	 \loigiai{
		Đường thẳng $d\colon x-y=3$  đi qua hai điểm $(0;-3)$, $(3;0)$.\\
		Xét điểm $O(0;0)\notin d$, ta có  $0-0=0<3$.\\
		Do đó $O(0;0)$  thuộc miền nghiệm của bất phương trình $x-y<3$.\\
		Vậy  phần không gạch sọc (chứa điểm $O$) trong hình dưới đây là miền nghiệm của bất phương trình $x-y<3$.
		\begin{center}
			\begin{tikzpicture}[line join=round,line cap=round,>=stealth,thick,scale=0.8]
					\tikzset{every node/.style={scale=0.9}} 
					\begin{scope}
						\clip (-1,-4) rectangle (4,1); 
						 \fill[pattern=north west lines] (-2,-5)--(5,-5)--(5,2)--cycle;
						\draw  (4,1)--(-1,-4);
					\end{scope}
					\draw[->] (-1,0)--(4,0) node[below]{$x$}; 
					\draw[->] (0,-4)--(0,1) node[left]{$y$}; 
					\draw  (0,0) node[below left]{$O$};
					\foreach \x in  {3}
					\draw[thin] (\x,1pt)--(\x,-1pt) node [above]  {$\x$};
					\foreach  \y in {-3}
					\draw[thin] (1pt,\y)--(-1pt,\y)  node [left] {$\y$};
				\end{tikzpicture}
		\end{center}
	 }
\end{ex}

% nguồn: MaHoaNH_TracNghiem_Lop10
\begin{ex}%[0D2H1-2]%[ma:12f7ec86]
	Biểu diễn  miền nghiệm của bất phương trình $2x-3y-6\le 0$ trên mặt phẳng tọa độ $Oxy$ là phần không gạch sọc trong hình nào sau đây?
	\choice 
		{\True 
		\begin{tikzpicture}[line join=round,line cap=round,>=stealth,thick]
				\tikzset{every  node/.style={scale=0.9}}
				\begin{scope}
					\clip (-2,-3)  rectangle (4,1.5);
					 \fill[pattern=north west lines] (-3,-4)--(6,-4)--(6,2)--cycle;
					 \draw (5.25,1.5)--(-1.5,-3);
				\end{scope}
				\draw[->]  (-2,0)--(4,0) node[below]{$x$};
				\draw[->] (0,-3)--(0,1.5)  node[left]{$y$};
				\draw (0,0) node[below  left]{$O$};
				\foreach \x in {-1,1,2,3} 
				\draw[thin] (\x,1pt)--(\x,-1pt) node [below]  {$\x$};
				 \foreach \y in {-2,-1,1}
				 \draw[thin] (1pt,\y)--(-1pt,\y) node [left] {$\y$};
			\end{tikzpicture}
		 }
		{ 
		\begin{tikzpicture}[line join=round,line cap=round,>=stealth,thick]
				 \tikzset{every node/.style={scale=0.9}}
				\begin{scope}
					\clip (-2,-1.5) rectangle (4,3); 
					\fill[pattern=north east lines] (-3,4)--(6,4)--(6,-2)--cycle; 
					 \draw (-1.5,3)--(5.25,-1.5);
				\end{scope}
				 \draw[->] (-2,0)--(4,0) node[below]{$x$};
				\draw[->] (0,-1.5)--(0,3) node[left]{$y$}; 
				\draw (0,0)  node[below left]{$O$};
				\foreach \x in  {-2,-1,1,2,3}
				\draw[thin] (\x,1pt)--(\x,-1pt)  node [below] {$\x$};
				\foreach \y  in {-1,1,2}
				\draw[thin]  (1pt,\y)--(-1pt,\y) node [left] {$\y$};
			\end{tikzpicture}
		} 
		 {
		\begin{tikzpicture}[line join=round,line cap=round,>=stealth,thick]
				\tikzset{every node/.style={scale=0.9}} 
				\begin{scope}
					 \clip (-2,-3) rectangle (4,1.5);
					\fill[pattern=north west  lines] (-3,-4)--(-3,2)--(6,2)--cycle;
					 \draw (5.25,1.5)--(-1.5,-3);
				\end{scope}
				\draw[->] (-2,0)--(4,0) node[below]{$x$}; 
				 \draw[->] (0,-3)--(0,1.5) node[left]{$y$};
				 \draw (0,0) node[below left]{$O$};
				 \foreach \x in {-1,1,2,3}
				\draw[thin] (\x,1pt)--(\x,-1pt) node [below] {$\x$}; 
				\foreach \y in {-2,-1,1} 
				 \draw[thin] (1pt,\y)--(-1pt,\y) node [left] {$\y$};
			\end{tikzpicture}
		 }
		 {
		\begin{tikzpicture}[line join=round,line cap=round,>=stealth,thick]
				\tikzset{every  node/.style={scale=0.9}}
				\begin{scope}
					 \clip (-2,-1.5) rectangle (4,3);
					\fill[pattern=north  east lines] (-3,4)--(-3,-2)--(6,-2)--cycle;
					\draw  (-1.5,3)--(5.25,-1.5);
				\end{scope}
				\draw[->]  (-2,0)--(4,0) node[below]{$x$};
				\draw[->] (0,-1.5)--(0,3)  node[left]{$y$};
				 \draw (0,0) node[below left]{$O$};
				\foreach \x  in {-1,1,2,3}
				\draw[thin] (\x,1pt)--(\x,-1pt) node [below] {$\x$}; 
				\foreach \y  in {-1,1,2}
				\draw[thin] (1pt,\y)--(-1pt,\y)  node [left] {$\y$};
			\end{tikzpicture}
		} 
	\loigiai{ 
		Đường thẳng $d\colon 2x-3y-6=0$ đi qua hai điểm  $(0;-2)$, $(3;0)$.\\
		Xét điểm $O(0;0)\notin d$, ta có $2\cdot0  -3\cdot0-6=-6<0$.\\
		Do đó $O(0;0)$ thuộc miền nghiệm của  bất phương trình $2x-3y-6\le 0$.\\
		Vậy phần không gạch sọc (chứa điểm $O$) trong hình dưới  đây là miền nghiệm của bất phương trình $2x-3y-6\le 0$.
		\begin{center}
			\begin{tikzpicture}[line join=round,line cap=round,>=stealth,thick]
					 \tikzset{every node/.style={scale=0.9}}
					\begin{scope}
						\clip (-2,-3) rectangle (4,1.5); 
						\fill[pattern=north  west lines] (-3,-4)--(6,-4)--(6,2)--cycle;
						 \draw (5.25,1.5)--(-1.5,-3);
					\end{scope}
					 \draw[->] (-2,0)--(4,0) node[below]{$x$};
					 \draw[->] (0,-3)--(0,1.5) node[left]{$y$};
					\draw (0,0)  node[below left]{$O$};
					\foreach  \x in {-1,1,2,3}
					\draw[thin] (\x,1pt)--(\x,-1pt) node [below]  {$\x$};
					\foreach \y in  {-2,-1,1}
					\draw[thin]  (1pt,\y)--(-1pt,\y) node [left] {$\y$};
				\end{tikzpicture}
		\end{center}
	 }
\end{ex}

% nguồn: MaHoaNH_TracNghiem_Lop10
\begin{ex}%[0D2Y1-2]%[ma:30226730]
	 Trong một buổi khiêu vũ có $20$ học sinh nam và $18$ học sinh nữ. Hỏi có bao nhiêu cách chọn ra một đôi nam nữ để khiêu vũ?
	 \choice
		 {$703$}
		{\True  $360$}
		 {$1406$}
		 {$38$}
	 \loigiai{
		Có $20$ cách chọn $1$ học sinh nam và $18$ cách chọn 1 học sinh nữ để tạo thành một cặp, nên có  $360$ cách chọn một đôi nam nữ để khiêu vũ.
	 }
\end{ex}

% nguồn: MaHoaNH_TracNghiem_Lop10
\begin{ex}%[0D2B1-2]%[ma:d5e64c85]
	Với năm chữ số $1$, $2$, $3$, $4$, $5$  có thể lập được bao nhiêu số tự nhiên có $2$ chữ số đôi một khác nhau?
	\choice 
		{\True $20$} 
		{$9$} 
		{$25$} 
		 {$10$}
	\loigiai{ 
		Với năm chữ số $1$, $2$, $3$, $4$, $5$, mỗi cách lấy $2$ chữ số khác nhau trong đó để sắp thứ tự chúng ta được một số  tự nhiên thỏa mãn yêu cầu bài toán.\\
		Do đó có $\mathrm{A}_5^2=20$ số thỏa mãn yêu cầu bài toán. 
	 }
\end{ex}

% nguồn: MaHoaNH_TracNghiem_Lop10
\begin{ex}%[0D2B1-2]%[ma:6978317f]
	Tập xác định của hàm số $y=\dfrac{x+2}{\sqrt{x-1}}+\sqrt{3-x}$ là 
	\choice 
		 {$[1;3]$}
		 {\True $(1;3]$}
		 {$(-\infty;3]$}
		 {$(1;+\infty)$}
	 \loigiai{
		$$\dfrac{x+2}{\sqrt{x-1}}+\sqrt{3-x}\text{ có  nghĩa }\Leftrightarrow \heva{&x-1>0\\&3-x\geq 0}\Leftrightarrow 1<x\leq 3.$$
		Vậy tập  xác định của hàm số đã cho là $(1;3]$.
	} 
\end{ex}

% nguồn: MaHoaNH_TracNghiem_Lop10
\begin{ex}%[0D2Y1-2]%[ma:982c560c]
	Từ thành phố A đến thành phố B có $6$ con đường, từ thành phố B đến thành phố C có $7$ con đường. Có bao nhiêu cách đi từ thành phố A đến thành phố C, biết phải đi  qua thành phố B?
	\choice 
		 {\True $42$}
		{$46$} 
		 {$48$}
		{$44$} 
	 \loigiai{
		\begin{itemize}
			\item Đi từ thành phố A đến thành phố B  có $6$ cách.
			\item Từ thành  phố B đến thành phố C có $7$ cách.
		\end{itemize}
		Theo quy tắc nhân ta có $7\cdot 6=42$ cách đi từ thành phố A đến thành phố C, biết phải đi  qua thành phố B.
	} 
\end{ex}

% nguồn: MaHoaNH_TracNghiem_Lop10
\begin{ex}%[0D2Y1-2]%[ma:e675919a]

	Mã xác thực (OTP - One Time Password) do một ngân hàng  gửi vào điện thoại của khách hàng cho mỗi lần giao là một dãy 6 kí tự từ các chữ số từ $0$ đến $9$. Có thể tạo ra bao nhiêu mã xác thực khác nhau như vậy?
	\begin{center}
		 \textbf{\texttt{Mã OTP xác thực giao dịch là 712892, hiệu lực trong 1 phút$\ldots$}}
	\end{center}
	 \choice
		{$151200$} 
		{$136080$} 
		{\True  $1000000$}
		{$900000$} 
	 \loigiai{
		 Mỗi kí tự trong dãy đều có $10$ cách chọn nên theo quy tắc nhân có thể tạo ra $10^6$ dãy.
	 }
\end{ex}

% nguồn: MaHoaNH_TracNghiem_Lop10
\begin{ex}%[0D2Y1-2]%[ma:39e7ffa1]
	Một người vào cửa hàng ăn, người đó chọn thực đơn gồm $1$ món ăn trong $5$ món, $1$ loại quả tráng miệng trong $5$ loại quả tráng miệng và một nước uống  trong $3$ loại nước uống. Có bao nhiêu cách chọn thực đơn?
	 \choice
		 {$25$}
		{\True $75$} 
		 {$100$}
		 {$15$}
	 \loigiai{
		\begin{itemize}
			\item Chọn $1$  món ăn có $5$ cách.
			\item Chọn $1$ quả tráng miệng có $5$ quả. 
			\item Chọn $1$ loại  nước uống có $3$ cách.
		\end{itemize}
		Theo quy tắc nhân ta có $5\cdot 5\cdot 3=75$ cách chọn thực  đơn.
	} 
\end{ex}

% nguồn: MaHoaNH_TracNghiem_Lop10
\begin{ex}%[0D2Y1-2]%[ma:4cd58a38]
	Cho tập $A=\left\{0;1;2;3;4;5\right\}$.  Có bao nhiêu số tự nhiên lẻ gồm $2$ chữ số được lập từ tập $A$?
	\choice 
		{$9$} 
		 {$12$}
		{\True  $15$}
		 {$18$}
	\loigiai{ 
		Gọi $\overline{ab}$ là số cần lập. 
		\begin{itemize}
			\item Chọn chữ số $b$ có $3$ cách  do $b\in\{1;3;5\}$.
			\item Chọn  chữ số $a$ có $5$ cách do $a\ne 0$.
		\end{itemize}
		Do đó, có  $5\cdot 3=15$ số được lập.
	 }
\end{ex}

% nguồn: MaHoaNH_TracNghiem_Lop10
\begin{ex}%[0D2B1-2]%[ma:c605b30c]
	Có tất cả bao nhiêu số  tự nhiên gồm $3$ chữ số?
	\choice 
		{$\mathrm{A}_{10}^{3}$} 
		 {\True $900$}
		{$\mathrm{C}_{ 10}^{3}$} 
		 {$\mathrm{C}_{9}^{3}$}
	 \loigiai{
		Gọi $n=\overline{abc}$ là  số cần lập.
		\begin{itemize}
			\item $a\ne 0$ nên $a$ có  $9$ cách chọn.
			\item $b$ có  $10$ cách chọn.
			\item $c$ có  $10$ cách chọn.
		\end{itemize}
		Vậy có  tất cả $9\cdot 10\cdot 10=900$ cách.
	} 
\end{ex}

% nguồn: MaHoaNH_TracNghiem_Lop10
\begin{ex}%[0D2G1-2]%[ma:56d07591]
	Tìm các giá trị thực của tham số $m$ để hàm số $y=\dfrac{x+m+2}{x-m}$ xác  định trên $(-1;2)$.
	\choice 
		{$\left\{\begin{aligned}
		&m\leq -1\\ 
		&m\geq 2\\ 
\end{aligned}\right. $}
{$\left[\begin{aligned}
&m<-1\\ 
&m>2\\ 
\end{aligned}\right. $ }
{$-1<m<2$} 
{\True $\left[\begin{aligned}
&m\leq  -1\\
 &m\geq 2\\
\end{aligned}\right. $ }
\loigiai{ 
	Hàm số  $y=\dfrac{x+m+2}{x-m}$ xác định khi $x\ne m$.\\
	Để hàm số $y=\dfrac{x+m+2}{x-m}$ xác định trên $(-1;2)$ khi và chỉ khi $\left[\begin{aligned}
	 &m\leq -1\\
	&m\geq 2.\\ 
\end{aligned}\right. $}
\end{ex}

% nguồn: MaHoaNH_TracNghiem_Lop10
\begin{ex}%[0D2B1-2]%[ma:dbcad538]
	 \immini{
	Nửa mặt phẳng không tô màu (kể cả bờ) trong hình vẽ bên  là miền nghiệm của bất phương trình nào?
	 \choice
		{$x+2y\ge  2$}
		{\True $2x+y\le  2$}
		 {$x+2y\le 2$}
		{$2x+y\ge 2$} 
		} 
		{ 
		\begin{tikzpicture}[line join=round,line cap=round,>=stealth,thick]
				 \tikzset{every node/.style={scale=0.9}}
				\begin{scope}
					 \clip (-1,-1) rectangle (2,3);
					\fill[pattern=crosshatch  dots gray] (-2,6)--(3,6)--(3,-4)--cycle;
					\draw  (-0.5,3)--(1.5,-1);
				\end{scope}
				\draw[->] (-1,0)--(2,0)  node[below]{$x$};
				\draw[->] (0,-1)--(0,3) node[left]{$y$}; 
				\draw  (0,0) node[below left]{$O$};
				\foreach \x in {1} 
				\draw[thin] (\x,1pt)--(\x,-1pt)  node [below] {$\x$};
				 \foreach \y in {2}
				\draw[thin] (1pt,\y)--(-1pt,\y)  node [left] {$\y$};
			\end{tikzpicture}
		} 
	 \loigiai{
		Đường thẳng $y=-2x+2$ đi qua điểm $A(1;0)$ và $B(0;2)$.\\ 
		Gốc tọa độ $O$ thuộc miền  nghiệm của bất phương trình nên hình vẽ chính là miền nghiệm của bất phương trình $2x+y\le 2$.
	 }
\end{ex}


%%%%%%%%%%%%%%%% ĐÚNG/SAI %%%%%%%%%%%%%%%%

% nguồn: MaHoaNH_DungSai_Lop10
\begin{ex}%[0D2N1-2]%[ma:64ca4026]
	Cho bất phương  trình $2x-1\le 0$ (1).
	 \choiceTF
		{ Điểm  $A(2;3)$ thuộc miền nghiệm của (1)}
		{\True Điểm $B(-10;100)$ thuộc miền  nghiệm của (1)}
		{ \True Điểm $C(5;2)$ không  thuộc miền nghiệm của (1)}
		{Điểm $D(-1;1)$ không thuộc miền nghiệm của  (1)}
	 \loigiai{
		\begin{itemchoice}
			\itemch Điểm  $A(2;3)$ không thuộc miền nghiệm của (1).
			\itemch Điểm $B(-10;100)$ thuộc  miền nghiệm của (1).
			\itemch Điểm $C(5;2)$ không thuộc  miền nghiệm của (1).
			\itemch  Điểm $D(-1;1)$ thuộc miền nghiệm của (1)
		\end{itemchoice}
	 }
\end{ex}

% nguồn: MaHoaNH_DungSai_Lop10
\begin{ex}%[0D2N1-2]%[ma:65c95007]
	Cho  bất phương trình $4x-3y\le 5\quad (*)$. Khi đó
	\choiceTF 
		{$(1;-1)$ là nghiệm của bất phương  trình $(*)$}
		{\True $(0;0)$ là  nghiệm của bất phương trình $(*)$}
		{\True $(2;1)$ là nghiệm của  bất phương trình $(*)$}
		{$(3;-1)$ là nghiệm của bất  phương trình $(*)$}
	\loigiai{ 
		\begin{itemchoice}
			\itemch  Thay $(-1;1)$ vào bất phương trình $(*)$$\Leftrightarrow 4\cdot 1-3\cdot (-1)\le 5 \Leftrightarrow 7\le 5$.
			\itemch  Thay $(0;0)$ vào bất phương trình $(*)$$\Leftrightarrow 4\cdot 0-3\cdot 0\le 5 \Leftrightarrow 0\le 5$.
			\itemch Thay $(2;1)$ vào bất phương trình  $(*)$$\Leftrightarrow 4\cdot 2-3\cdot 1\le 5 \Leftrightarrow 3\le 5$.
			\itemch Thay $(3;-1)$ vào bất phương trình  $(*)$$\Leftrightarrow 4\cdot 3-3\cdot (-1)\le 5 \Leftrightarrow 15\le 5$.
		\end{itemchoice}
	} 
\end{ex}

% nguồn: MaHoaNH_DungSai_Lop10
\begin{ex}%[0D2N1-2]%[ma:726f60d1]
	Cho bất phương trình  $2(x+1)+2(y+3)\le 12$ (1).
	 \choiceTF
		{ Điểm  $A(-1;1)$ không thuộc miền nghiệm của (1)}
		{ Điểm $B(0;3)$  thuộc miền nghiệm của (1)}
		{  \True Điểm $C\left(\dfrac{5}{2};0\right)$ thuộc miền nghiệm của (1)}
		{Phần không bị gạch (kể cả bờ) trong  hình vẽ bên dưới biểu diễn miền nghiệm của (1).
		\begin{tikzpicture}[scale=0.5,font=\footnotesize,line join=round,line cap=round,>=stealth]
				 \draw (-2,4)--(4,-2);
				\begin{scope}
					\clip (-2,-2) rectangle  (4,4);
					\fill[pattern=north  east lines] (-2.1,4.1)--(-2.1,-2.1)--(4.1,-2.1)--cycle;
				\end{scope}
				\draw[->]  (-2,0)--(4,0) node[right]{$x$};
				 \draw[->] (0,-2)--(0,4) node[above]{$y$};
				\draw[fill=black] (0,0)  circle (1pt) node[above left=-2pt] {$O$};
				 \draw[fill=black] (0,2) circle (1pt) node[right]{$2$}
					(2,0)  circle (1.1pt) node[above]{$2$};
				\clip  (-2,-1) rectangle (2.5,4.5);
			\end{tikzpicture}
		 }
	\loigiai{Ta có $2(x+1)+2(y+3)\le  12 \Leftrightarrow x+y\le 2$.
		\begin{itemchoice}
			\itemch Điểm $A(-1;1)$ thuộc miền nghiệm của (1). 
			 \itemch Điểm $B(0;3)$ không thuộc miền nghiệm của (1).
			\itemch Điểm $C\left(\dfrac{5}{2};0\right)$ thuộc  miền nghiệm của (1).
			 \itemch Miền nghiệm trong hình vẽ không chứa điểm $O$ nên không phải là miền nghiệm của (1).
		\end{itemchoice}
	 }
\end{ex}

% nguồn: MaHoaNH_DungSai_Lop10
\begin{ex}%[0D2N1-2]%[ma:89919cbf]
	 \immini
	{ 
	Phần không bị gạch (không kể bờ) trong hình vẽ bên biểu diễn miền nghiệm của một bất phương trình  bậc nhất hai ẩn.
	 }
	{ 
	\begin{tikzpicture}[>=stealth,line join=round,line cap=round,font=\footnotesize,scale=1]
			\begin{scope}[yscale=.5]
				\draw[->](-1,0)--(4,0) node[below]{$x$}; 
				\draw[->](0,-4)--(0,2) node[right]{$y$}; 
				\draw[domain=-.5:2.5]  plot (\x,{2*\x-3});
				\fill (0,-3) circle (0.04)  node[left]{$-3$}
					(0,0) circle (0.04) node[below  left]{$O$}
					(1.5,0) circle (0.04) node[above left]{$\dfrac{3}{2}$}; 
				\fill[pattern=north  east lines](-.5,-4)--(2.5,2)--(4,2)--(4,-4)--cycle;
			\end{scope}
		\end{tikzpicture}
	 }
	 \choiceTF
		{\True Điểm $A(-1;3)$ thuộc miền nghiệm  của bất phương trình đã cho}
		{\True Điểm $B(0;-3)$ không thuộc miền nghiệm của bất phương trình  đã cho}
		{ Điểm $O(0;0)$ không thuộc miền nghiệm của  bất phương trình đã cho}
		 {Miền nghiệm trong hình vẽ là miền nghiệm của bất phương trình $2x-y\le 3$ }
	 \loigiai{
		\begin{itemchoice}
			\itemch Điểm $A(-1;3)$ thuộc miền nghiệm của bất phương  trình đã cho.
			\itemch Điểm $B(0;-3)$ không thuộc miền nghiệm của bất phương trình  đã cho.
			\itemch Điểm $O(0;0)$ thuộc miền nghiệm của bất phương trình  đã cho.
			\itemch  Miền nghiệm trong hình vẽ là miền miền nghiệm của bất phương trình $2x-y<3$.
		\end{itemchoice}
	} 
\end{ex}

% nguồn: MaHoaNH_DungSai_Lop10
\begin{ex}%[0D2H1-2]%[ma:4ecb1cad]
	Xét tính đúng, sai của các mệnh đề sau: 
	 \choiceTF
		{\True Cho  bất phương trình $3-2y> 0$ có miền nghiệm là nửa mặt phẳng bờ 3$ -2y=0$ chứa $O$ (bỏ bờ)}
		{Cho bất phương trình  $2x+y > 1$ có miền nghiệm là nửa mặt phẳng bờ $-2x+y+1=0$ chứa $O$ (bỏ bờ)}
		{Cho bất phương trình $-2x+y+1\leq 0$ có miền nghiệm là nửa mặt phẳng  bờ $-2x+y+1=0$ chứa $O$}
		{\True Cho bất phương trình $2x-3y+5\geq 0$ có miền  nghiệm là nửa mặt phẳng bờ $2x-3y+5=0$ chứa $O$}
	 \loigiai{
		\begin{itemchoice}
			\itemch Điểm $O(0; 0)$ có tọa độ thỏa mãn bất phương trình, do đó miền nghiệm là nửa mặt phẳng bờ 3$ -2y=0$ chứa  $O$ (bỏ bờ).
			\itemch Miền nghiệm là nửa mặt phẳng bờ $-2x+y+1=0$ không chứa $O$  (bỏ bờ).
			\itemch Miền nghiệm là nửa mặt phẳng  bờ $-2x+y+1=0$ không chứa $O$.
			\itemch Miền nghiệm là nửa  mặt phẳng bờ $2x-3y+5=0$ chứa $O$.
		\end{itemchoice}
	} 
\end{ex}

% nguồn: MaHoaNH_DungSai_Lop10
\begin{ex}%[0D2H1-2]%[ma:d67e2a60]
	Cho bất phương  trình $3 x-5 y\le-15$. Khi đó
	 \choiceTF
		{Điểm $O(0;0)$ thuộc  miền nghiệm của bất phương trình đã cho}
		{\True Đường thẳng $d\colon 3x -5y =-15$  đi qua hai điểm $A(-5; 0)$ và $B(0; 3)$}
		{Miền nghiệm của bất phương trình là nửa mặt  phẳng bờ $d\colon 3x -5y =-15$ (kể cả đường thẳng $d$) chứa điểm $(O)$}
		{\True Có vô số cặp số $(x; y)$ sao cho $x<0$, $y>0$, $x$, $y\in \mathbb{Z}$  là nghiệm của bất phương trình đã cho}
	 \loigiai{
		\begin{itemchoice}
			\itemch  Vì $3\cdot0-5\cdot0=0>-15$.\\
			Do đó điểm $O(0;0)$ không thuộc miền nghiệm của bất phương trình $3  x-5 y\le-15$.
			\itemch Vì $3\cdot (-5)-5\cdot0=-15 \Rightarrow A(-5; 0)\in  d$ và $3\cdot 0-5\cdot3=-15 \Rightarrow B(0; 3)\in d$.\\
			Vậy đường thẳng $d\colon 3x -5y =-15$ đi qua hai điểm $A(-5; 0)$ và  $B(0; 3)$.
			\itemch Vì $O\notin d$ và  $3\cdot0-5\cdot0=0>-15$.\\
			Do đó miền nghiệm của bất phương trình là nửa mặt phẳng bờ $d\colon 3x -5y =-15$ (kể cả đường thẳng $d$) không chứa điểm $(O)$. 
			\itemch  Ta có $3 x-5 y\le-15\Leftrightarrow 3x\le 5y -15\Leftrightarrow x\le \dfrac{5}{3}y-3$.\\
			Vì $x<0$, $y>0$, $x$ $y\in \mathbb{Z}$ nên $y\in\{3; 6; 9; \ldots\}$.\\ 
			Vậy có vô số cặp số  $(x; y)$ sao cho $x<0$, $y>0$, $x$, $y\in \mathbb{Z}$ là nghiệm của bất phương trình đã cho.
		\end{itemchoice}
	} 
\end{ex}

% nguồn: MaHoaNH_DungSai_Lop10
\begin{ex}%[0D2H1-2]%[ma:43fa4ece]
	 Cho bất phương trình bậc nhất hai ẩn: $x-2y+2\leq 0$. Khi đó
	\choiceTF 
		{\True Miền nghiệm của bất phương trình $x-2y+2\leq 0$ là nửa  mặt phẳng kể cả bờ $x-2y+2=0$, không chứa gốc tọa độ $O$}
		 {\True $(1; 4)$ là nghiệm của bất phương trình $x-2y+2\leq 0$}
		{$(0; 3)$ không là nghiệm  của bất phương trình $x-2y+2\leq 0$}
		{$(2; 2)$ không là nghiệm của bất phương trình $x-2y+2\leq 0$} 
	\loigiai{ 
		\begin{itemchoice}
			\itemch Vẽ đường thẳng $d\colon x-2y+2=0$ đi qua hai điểm $A(0; 1)$ và  $B(-2; 0)$.
		 \immini{
			 Xét gốc tọa độ $O(0; 0)$ ta thấy $O\notin d$ và $1\cdot 0-2\cdot 0+2=2> 0$.\\
			Do đó, miền nghiệm của bất phương trình $x-2y+2\leq 0$ là nửa mặt phẳng kể cả bờ $d$ không chứa gốc tọa độ $O$ (miền không  bị tô đậm trong hình).}
			{\begin{tikzpicture}[scale=0.9,font=\footnotesize,line join=round,line cap=round,>=stealth]
		\draw[->] (-3,0)--(4,0)  node[below] {$x$};
		\draw[->]  (0,-1)--(0,3) node[left] {$y$};
		 \fill (0,0) circle (1pt) node[below left]{$O$};

		\fill 
			(0,1) circle (1pt)  node[left]{$A$}
			(-2,0)  circle (1pt) node[above]{$B$};
		 \fill[pattern=north east lines,smooth] plot[domain=-3:4](\x,{1/2*(\x)+1})-- plot[domain=4:-3](\x,{0*\x-1}) -- cycle;
		\draw[thick,smooth,samples=200]  plot[domain=-3:4] (\x,{1/2*(\x)+1});
	\end{tikzpicture}}
			\itemch Ta có $1\cdot 1-2\cdot  4+2=-5< 0$ nên $(1; 4)$ là nghiệm của bất phương trình $x-2y+2\leq 0$.
			\itemch Ta có $1\cdot  0-2\cdot 3+2=-3< 0$ nên $(0; 3)$ là nghiệm của bất phương trình $x-2y+2\leq 0$.
			\itemch Ta có $1\cdot 2-2\cdot  2+2=0$ nên $(2; 2)$ là nghiệm của bất phương trình $x-2y+2\leq 0$.
		\end{itemchoice}
	} 
\end{ex}


%%%%%%%%%%%%%%%% TRẢ LỜI NGẮN %%%%%%%%%%%%%%%%

% nguồn: MaHoaNH_TraLoiNgan_Lop10
\begin{bt}%[0D2H1-2]%[ma:3c186023]
	 Tìm các giá trị nhỏ nhất của tham số $m$ sao cho $\heva{&x=1\\
	&y=-1}$ là nghiệm của bất phương trình $m \dfrac{x}{2}-(m+1) y+2\geq 0$. 
	 \shortans{-2}
	 \loigiai{
		Ta  có $\heva{&x=1\\
		&y=-1}$ là nghiệm của bất phương trình $m \dfrac{x}{2}-(m+1) y+2\geq 0$ khi và chỉ khi \[m \dfrac{1}{2}-(m+1)(-1)+2\geq 0\Leftrightarrow \dfrac{3}{2} m+3\geq 0\Leftrightarrow m \geq-2.\] 
	} 
\end{bt}

% nguồn: MaHoaNH_TraLoiNgan_Lop10
\begin{bt}%[0D2H1-2]%[ma:02b03b70]
	Có tất cả bao nhiêu giá trị nguyên của tham số $m$ thuộc đoạn $[-2025;2025]$ để nghiệm của hệ phương trình $\heva{&2x+y=5\\&x-3y=-1}$ không thuộc miền nghiệm của bất phương trình  $-x+(m+1)y+3\ge 0$?
	\shortans{2023} 
	 \loigiai{
		Ta  có $\heva{&2x+y=5\\&x-3y=-1} \Leftrightarrow \heva{&2x+y=5\\&2x-6y=-2} \Leftrightarrow \heva{&2x+y=5\\&7y=7} \Leftrightarrow \heva{&2x+y=5\\&y=1} \Leftrightarrow \heva{&x=2\\&y=1.}$\\
		Để $(2;1)$ không thuộc miền nghiệm của bất phương trình $-x+(m+1)y+3\ge 0$ thì  nghiệm $(2;1)$ là nghiệm của bất phương trình $-x+(m+1)y+3<0$.\\
		 Khi đó ta có $-2+m+1+3< 0\Leftrightarrow m+2< 0\Leftrightarrow m <-2$.\\
		Mà $m$ nguyên, $m\in  \left[-2025;2025\right]$ nên $m\in \left\{-2025;\ldots;-3\right\}$.\\
		Vậy có $2023$ giá trị nguyên  $m$ thỏa mãn.
	 }
\end{bt}

% nguồn: MaHoaNH_TraLoiNgan_Lop10
\begin{bt}%[0D2H1-2]%[ma:96bb41ee]
	Cho bất phương trình $\dfrac{x}{2}+\dfrac{y}{3}-1\leq 0$. Miền nghiệm có chứa bao nhiêu điểm $(x;  y)$ với $x, y$ là các số nguyên dương?
	 \shortans{1}
	 \loigiai{
		Do $x > 0$, $\dfrac{x}{2}+\dfrac{y}{3}-1\leq 0$ nên ta có $\dfrac{y}{3} < 1\Leftrightarrow  y < 3$.\\
		Do $y$ nguyên dương nên $y  \in\{1; 2\}$.\\
		 Với $y=1$, ta có $\heva{&\dfrac{x}{2}+\dfrac{1}{3}-1\leq 0\\
		&x > 0} \Leftrightarrow 0< x  \leq \dfrac{4}{3} \Leftrightarrow x=1$.\\
		Với $y=2$, ta có $\heva{&\dfrac{x}{2}+\dfrac{2}{3}-1\leq 0\\ 
		&x > 0} \Leftrightarrow 0< x \leq \dfrac{2}{3} \Leftrightarrow x \in  \varnothing$.\\
		Vậy bất phương trình $\dfrac{x}{2}+\dfrac{y}{3}-1\leq 0$ có  nghiệm nguyên dương là $(1; 1)$.
	 }
\end{bt}

% nguồn: MaHoaNH_TraLoiNgan_Lop10
\begin{bt}%[0D2H1-2]%[ma:44c00508]
	Cho tam giác $ABC$ có $A(0; 3)$; $B(-1; 2)$; $C(2; 1)$. Số giá trị nguyên của tham số $m$ để điểm $M\left(m; \dfrac{2m-1}{2}\right)$ nằm bên trong tam giác  $ABC$ là bao nhiêu?
	 \shortans{0}
	\loigiai{ 
		 Đường thẳng $AB\colon \dfrac{x-0}{-1-0}=\dfrac{y-3}{2-3} \Leftrightarrow x-y+3=0$.\\
		Đường  thẳng $AC\colon \dfrac{x-0}{2-0}=\dfrac{y-3}{1-3} \Leftrightarrow x+y-3=0$.\\
		Đường thẳng $BC\colon \dfrac{x-2}{2-(-1)}=\dfrac{y-1}{1-2} \Leftrightarrow  x+3y-5=0$.\\
		Điều kiện cần và đủ để điểm $M$ nằm bên  trong tam giác $ABC$ là điểm $M$ cùng với mỗi đỉnh $A$, $B$, $C$ lần lượt cùng phía với nhau đối với cạnh $AB$, $AC$, $BC$.\\
		Điều này xảy  ra khi và chỉ khi
		 \[\heva{&(1\cdot 0+3\cdot 3-5) \cdot\left(1\cdot m+3\cdot \dfrac{2m-1}{2}-5\right) > 0\\
		&(1\cdot(-1)+1\cdot 2-3) \cdot\left(1\cdot  m+1\cdot \dfrac{2m-1}{2}-3\right) > 0\Leftrightarrow \\
		&(1\cdot 2-1\cdot 1+3) \cdot\left(1\cdot m-1\cdot  \dfrac{2m-1}{2}+3\right) > 0}\heva{&m > \dfrac{13}{8} \\
		&m <  \dfrac{7}{4} \\
		&14> 0\,(\text{thỏa mãn})} \Leftrightarrow \dfrac{13}{8} < m <  \dfrac{7}{4}.\]
	 }
\end{bt}

% nguồn: MaHoaNH_TraLoiNgan_Lop10
\begin{bt}%[0D2H1-2]%[ma:80a69078]
	\immini{Nửa mặt phẳng không bị gạch (không kể $d$) ở hình  bên là miền nghiệm của bất phương trình có dạng $ax+by-2>0$. Giá trị của $a+b$ bằng bao nhiêu?}
	{ 
	\begin{tikzpicture}[scale=1,line join=round,line cap=round,>=stealth]
			 \tikzset{every node/.style={scale=0.9}}
			\begin{scope}
				\clip (-1,-2.5) rectangle  (3.5,1.5);
				\fill[pattern=north  west lines] (-2,-4)--(-2,2.5)--(4.5,2.5)--cycle;
				\draw  (3.5,1.5)--(-0.5,-2.5);
			\end{scope}
			 \draw[->] (-1,0)--(3.5,0) node[below]{$x$};
			\draw[->] (0,-2.5)--(0,1.7)  node[left]{$y$};
			\draw  (0,0) node[below left]{$O$};
			 \foreach \x in {2,3}
			\draw[thin] (\x,1pt)--(\x,-1pt)  node [above] {$\x$};
			\foreach  \y in {-2,-1}
			\draw[thin] (1pt,\y)--(-1pt,\y) node [left] {$\y$}; 
		\end{tikzpicture}
	 }
	 \shortans{0}
	 \loigiai{
		Dựa vào hình ta thấy đường thẳng $d$ có  hệ số góc khác $0$ nên gọi $(d)\colon y=ax+b$ .\\
		Ta  có : $(2;0)\in (d)\Rightarrow 2a+b=0$ và $(0;-2)\in (d)\Rightarrow b=-2$,\\
		Vậy $(d)\colon y=x-2$ hay  $x-y-2=0$.\\
		Mà $(0;0)$ không thuộc miền nghiệm  $\Rightarrow x-y-2>0$.\\
		Do đó $a=1; b=-1  \Rightarrow a+b=0$.
	} 
\end{bt}

% nguồn: MaHoaNH_TraLoiNgan_Lop10
\begin{bt}%[0D2V1-2]%[ma:0ab02aaf]
	Cho bất phương trình  $2x-(m+2)y+m\le 0$. Có bao nhiêu giá trị nguyên âm của tham số $m$ để bất phương trình có một nghiệm là $(1;3)$?
	\shortans{2} 
	\loigiai{ 
		Thay  $(1;3)$ vào bất phương trình ta được $2\cdot 1-(m+2)\cdot 3+m\le 0\Leftrightarrow m\ge -2$.\\
		 Mà $m$ nguyên âm $\Rightarrow m=\left\lbrace -2;-1\right\rbrace$.\\
		Vậy có $2$  giá trị nguyên âm của tham số $m$ thỏa mãn bất phương trình đã cho.
	 }
\end{bt}

% nguồn: MaHoaNH_TraLoiNgan_Lop10
\begin{bt}%[0D2V1-2]%[ma:a5bb53ec]
	Tìm các giá trị nguyên của tham số $m\in [-10;10]$ sao cho $(1; -1)$ là nghiệm của bất phương trình $m\cdot\dfrac{x}{2}-(m+1) y+2\ge  0$.
	\shortans{13} 
	\loigiai{ 
		Ta có $(1; -1)$ là nghiệm của bất  phương trình $m\cdot\dfrac{x}{2}-(m+1) y+2\ge 0$.\\
		Suy ra $m\cdot\dfrac{1}{2}-(m+1)\cdot (-1)+2\ge 0\Leftrightarrow  \dfrac{3}{2}m+3\ge 0\Leftrightarrow m\ge -2$.\\
		Do  $m\in\mathbb{Z}$ và $m\in [-10;10]$ nên $m\in\left\lbrace -2;-1;0;1;\ldots;10 \right\rbrace $.\\
		Vậy có $13$ số nguyên $m$  thỏa mãn đề bài.
	 }
\end{bt}

% nguồn: MaHoaNH_TraLoiNgan_Lop10
\begin{bt}%[0D2V1-2]%[ma:40269ca1]
	Có bao nhiêu nghiệm $(x,y)$ của bất phương trình $\dfrac{x}{2}+\dfrac{y}{3}-1\le 0$  thỏa mãn $x,y$ là các số nguyên dương?
	\shortans{1} 
	\loigiai{ 
		Ta có $\dfrac{x}{2}+\dfrac{y}{3}-1\le  0 \Leftrightarrow
		\heva{&\dfrac{x}{2}\le 1-\dfrac{y}{3}\\&\dfrac{y}{3}\le 1- \dfrac{x}{2}.  }$\\
		Vì $x$, $y$  nguyên dương nên ta có $\heva{&\dfrac{y}{3}<1\\&\dfrac{x}{2}<1}
		 \Leftrightarrow \heva{&x<2\\&y<3}$.
		Suy  ra $x\in\{1\}$, $y\in\{1;2\}$.\\
		Với  $x=1\Rightarrow \dfrac{1}{2}+\dfrac{y}{3}\le 1\Leftrightarrow \dfrac{y}{3}\le \dfrac{1}{2}\Leftrightarrow y\le \dfrac{3}{2}\Rightarrow y=1$.\\
		Vậy bất phương trình $\dfrac{x}{2}+\dfrac{y}{3}-1\le 0$ có nghiệm nguyên dương là  $(1;1)$.
	 }
\end{bt}

% nguồn: MaHoaNH_TraLoiNgan_Lop10
\begin{bt}%[0D2V1-2]%[ma:b47ce965]
	Cho bất phương trình $x+2y\ge -4$. Miền  nghiệm có chứa bao nhiêu điểm $(x;y)$ với $x,y$ là các số nguyên âm?
	\shortans{2} 
	 \loigiai{
	\immini{Biểu diễn miền nghiệm của bất  phương trình $x+2y\ge -4$ ta được hình bên.\\
		 Dựa vào miền nghiệm trên hình ta thấy chỉ có $2$ điểm có hoành độ và tung độ đều nguyên âm là $(-1;-1)$ và $(-2;-1)$.
		 }
		{ 
		\begin{tikzpicture}[scale=1,line join=round,line cap=round,>=stealth]
				\draw[ultra thin,color=gray!50] (-5,-3) grid  (2,2);
				\tikzset{every  node/.style={scale=0.9}}
				\begin{scope}
					\clip (-5,-3) rectangle  (2,2);
					\fill[pattern=north east lines] (-9,2.5)--(-9,-3.5)--(3,-3.5)--cycle; 
					\draw  (-8,2)--(2,-3) node [pos=0.65,above,sloped] {$x+2y+4=0$};
				\end{scope}
				\draw[->] (-5,0)--(2,0)  node[above]{$x$};
				\draw[->] (0,-3)--(0,2)  node[right]{$y$};
				\draw (0,0) node[below  left]{$O$};
				 \foreach \x in {-1,-2,-3,-4}
				\draw[thin] (\x,1pt)--(\x,-1pt) node [above] {$\x$}; 
				\foreach \y  in {-1,-2}
				\draw[thin] (1pt,\y)--(-1pt,\y) node [right]  {$\y$};
			\end{tikzpicture}
		} 
	} 
\end{bt}

% nguồn: MaHoaNH_TraLoiNgan_Lop10
\begin{bt}%[0D2V1-2]%[ma:bd53aec7]
	Miền  nghiệm của bất phương trình $x+y\le 4$ cùng với hai trục tọa độ $Ox$, $Oy$ tạo thành một tam giác. Diện tích của tam giác đó bằng bao nhiêu?
	 \shortans{8}
	\loigiai{ 
	 \immini{
		Vẽ đường thẳng $d\colon x+y=4$ đi qua hai  điểm $(4;0)$ và $(0;4)$.\\
		Xét  điểm $O(0;0)\notin d$, ta có $0 + 0 <4$.\\
		Do đó điểm $O$ thuộc miển nghiệm của bất phương trình  $x+y\le 4$.\\
		Biểu diễn miền nghiệm của bất phương trình $x+y\le 4$  ta được hình bên.\\
		Diện tích  tam giác $S=\dfrac{1}{2}\cdot 4\cdot 4=8$.
		 }
		{ 
		\begin{tikzpicture}[scale=1,line join=round,line cap=round,>=stealth,scale=.8]
				 \tikzset{every node/.style={scale=0.8}}
				\begin{scope}
					\clip (-1,-1) rectangle  (5,5);
					\fill[pattern=north east  lines] (-2,6)--(6,6)--(6,-2)--cycle;
					\draw (-1,5)--(5,-1)  node [pos=0.45,below,sloped] {$d\colon x+y-4=0$};
				\end{scope}
				 \draw[->] (-1.5,0)--(5.5,0) node[below]{$x$};
				 \draw[->] (0,-1)--(0,5.2) node[left]{$y$};
				\draw (0,0)  node[below left]{$O$};
				\foreach  \x in {1,2,3,4}
				\draw[thin] (\x,1pt)--(\x,-1pt)  node [below] {$\x$};
				\foreach \y in {1,2,3,4} 
				\draw[thin] (1pt,\y)--(-1pt,\y) node [left] {$\y$}; 
			\end{tikzpicture}
		} 
	} 
\end{bt}


%%%%%%%%%%%%%%%% TỰ LUẬN %%%%%%%%%%%%%%%%

% nguồn: MaHoaNH_TuLuan_Lop10
\begin{bt}%[0D2H1-2]%[ma:067cccc4]
	Biểu diễn miền nghiệm của bất phương trình $\dfrac{5x+y-1}{2}\ge x+y+1$ trêm mặt  phẳng tọa độ $Oxy$.
	 \loigiai{
		Ta có $\dfrac{5x+y-1}{2}>x+y+1\Leftrightarrow 5x+y-1  >2x+2y+2\Leftrightarrow 3x-y-3\ge 0$.
		\begin{itemize}
			\item Vẽ đường thẳng $\Delta\colon 3x-y-3=0$ đi qua  hai điểm $\left(1;0\right)$ và $\left(0;-3\right)$.
			\item Xét điểm $O(0;  0)\notin \Delta$. Ta có $3\cdot0-0-3=-3<0$.
			\item Do đó, miền nghiệm của bất phương trình $3x-y-3\ge  0$ là nửa mặt phẳng bờ $\Delta$ chứa điểm $O$ kể cả đường thẳng $\Delta$.
		\end{itemize}
		\begin{center}
			\begin{tikzpicture}[scale=.9,>=stealth,font=\footnotesize,line join=round,line cap=round]
					 \def\a{3} \def\b{-3}
					 \def\xmin{-3} \def\xmax{3}
					\def\ymin{-4}  \def\ymax{2}
					\draw[->] (\xmin,0)--(\xmax+.2,0) node [below]{$x$}; 
					 \draw[->] (0,\ymin)--(0,\ymax+.2) node [left]{$y$};
					 \node at (0,0) [below left]{$O$};
					\clip (\xmin,\ymin) rectangle (\xmax,\ymax); 
					\fill[pattern=north east lines]  (\xmin,\ymin)--plot (\x,{\a*(\x)+\b})--(\xmin,\ymax) --cycle;
					\draw[smooth,samples=300]  plot (\x,{\a*(\x)+\b});
					\node[rotate=70,right] at  (.5,-3.2) {$\Delta\colon 3x -y -3= 0$};
					\draw[fill=black] (0,-3) circle (1pt) node[below right]{$-3$}; 
					\draw[fill=black] (1,0) circle (1pt) node[above right]{$1$}; 
				\end{tikzpicture}
		\end{center}
	 }
\end{bt}

% nguồn: MaHoaNH_TuLuan_Lop10
\begin{bt}%[0D2H1-2]%[ma:2f79cc25]
	Biểu diễn miền nghiệm của mỗi bất  phương trình $\dfrac{x+y}{2}>\dfrac{2x-y+1}{3}$ trên mặt phẳng tọa độ $Oxy$.
	\loigiai{ 
		 Ta có $\dfrac{x+y}{2}>\dfrac{2x-y+1}{3}\Leftrightarrow 3x+3y >4x-2y+2\Leftrightarrow x-5y+2<0$.
		\begin{itemize}
			\item Vẽ đường  thẳng $\Delta\colon x-5y+2=0$ đi qua hai điểm $\left(-2;0\right)$ và $\left(0;\dfrac{2}{5}\right)$.
			\item Xét điểm $O(0; 0)\notin \Delta$.  Ta có $0-5\cdot0+2=2>0$.
			\item Do đó, miền nghiệm của bất phương trình $x-5y+2<0$ là nửa mặt  phẳng bờ $\Delta$ không chứa điểm $O$ không kể đường thẳng $\Delta$.
		\end{itemize}
		\begin{center}
			\begin{tikzpicture}[scale=.9,>=stealth,font=\footnotesize,line join=round,line cap=round]
					\def\a{1/5} \def\b{2/5} 
					\def\xmin{-3.5} \def\xmax{4} 
					\def\ymin{-2.5}  \def\ymax{2}
					\draw[->]  (\xmin,0)--(\xmax+.2,0) node [below]{$x$};
					\draw[->]  (0,\ymin)--(0,\ymax+.2) node [left]{$y$};
					\node  at (0,0) [below left]{$O$};
					\clip (\xmin,\ymin)  rectangle (\xmax,\ymax);
					\fill[pattern=north east  lines] plot (\x,{\a*(\x)+\b})--(\xmax,\ymin)--(\xmin,\ymin) --cycle;
					\draw[smooth,samples=300] plot (\x,{\a*(\x)+\b}); 
					\node[rotate=12,right] at  (.5,.7) {$\Delta\colon x -5y +2= 0$};
					 \draw[fill=black] (0,.4) circle (1pt) node[above left]{$\dfrac{2}{5}$};
					\draw[fill=black] (-2,0) circle  (1pt) node[above]{$-2$};
				\end{tikzpicture}
		\end{center}
	} 
\end{bt}

% nguồn: MaHoaNH_TuLuan_Lop10
\begin{bt}%[0D2H1-2]%[ma:75fa2c13]
	Biểu diễn miền nghiệm của mỗi bất phương trình sau trên mặt phẳng toạ độ. 
	\begin{multicols}{2}
		\begin{enumerate}
			\item $3x +2y \ge  300$.
			\item $7x +20y<0$. 
			 \item $2x+3y-1<0$
			\item $x+3y-2\le 0$. 
			\item $\dfrac{x+y}{2}>\dfrac{2x-y+1}{3}$. 
			 \item $\dfrac{5x+y-1}{2}\ge x+y+1$.
		\end{enumerate}
	\end{multicols}
	 \loigiai{
		\begin{enumerate}
			\item Biểu diễn miền nghiệm của  mỗi bất phương trình $3x +2y \ge 300$.
			\begin{itemize}
				\item Vẽ đường thẳng $\Delta\colon 3x +2y =  300$ đi qua hai điểm $(100; 0)$ và $(0; 150)$.
				\item Xét điểm $O(0;  0)\notin \Delta$. Ta có $3\cdot0+2\cdot0=0<300$.
				\item Do đó, miền nghiệm của bất  phương trình $3x +2y \ge 300$ là nửa mặt phẳng bờ $\Delta$ không chứa điểm $O$ kể cả đường thẳng $\Delta$.
			\end{itemize}
			\begin{center}
				\begin{tikzpicture}[scale=1,>=stealth,font=\footnotesize,line join=round,line cap=round]
						\def\a{-3/2}  \def\b{3}
						\def\xmin{-1} \def\xmax{4} 
						\def\ymin{-1}  \def\ymax{4}
						\draw[->] (\xmin,0)--(\xmax,0) node  [below]{$x$};
						\draw[->] (0,\ymin)--(0,\ymax) node  [left]{$y$};
						\fill  (0,0) circle (1pt) node[below left]{$O$};
						 \clip (\xmin+0.1,\ymin+0.1) rectangle (\xmax-0.5,\ymax-0.1);
						\fill[pattern=north east  lines] (\xmin,\ymin) --(\xmin,\ymax)--plot (\x,{-1.5*(\x)+3})--cycle;
						\draw[smooth,samples=300] plot (\x,{\a*(\x)+\b}); 
						\node[rotate=-58,right] at (.5,3) {$\Delta\colon 3x +2y = 300$}; 
						\draw[fill=black]  (2,0) circle (1pt) node[above right]{$100$};
						\draw[fill=black] (0,3) circle (1pt)  node[above right]{$150$};
					\end{tikzpicture}
			\end{center}
			\item Biểu diễn miền nghiệm của mỗi bất phương trình  $7x +20y<0$.
			\begin{itemize}
				\item Vẽ đường thẳng $\Delta\colon 7x +20y=0$ đi qua hai điểm $(0; 0)$  và $\left(2;-\dfrac{7}{10}\right)$.
				\item Xét điểm $M(1; 1)\notin  \Delta$. Ta có $7\cdot1+20\cdot1=27>0$.
				 \item Do đó, miền nghiệm của bất phương trình $7x +20y<0$ là nửa mặt phẳng bờ $\Delta$ không chứa điểm $M$ không kể đường thẳng $\Delta$.
			\end{itemize}
			\begin{center}
				\begin{tikzpicture}[scale=1,>=stealth,font=\footnotesize,line join=round,line cap=round]
						\def\a{-7/20} \def\b{0} 
						\def\xmin{-3} \def\xmax{4} 
						\def\ymin{-2}  \def\ymax{2}
						\draw[->] (\xmin,0)--(\xmax+.2,0) node  [below]{$x$};
						 \draw[->] (0,\ymin)--(0,\ymax+.2) node [left]{$y$};
						\fill (0,0) circle (1pt)  node[below left]{$O$};
						\clip (\xmin,\ymin)  rectangle (\xmax,\ymax);
						\fill[pattern=north east  lines] plot (\x,{\a*(\x)+\b})--(\xmax,\ymax) --(\xmin,\ymax)--cycle;
						\draw[smooth,samples=300] plot (\x,{\a*(\x)+\b}); 
						\node[rotate=-20,right] at (.5,.-.5)  {$\Delta\colon 7x +20y = 0$};
						\draw[fill=black] (1,1) circle (1pt) node[above right,fill=white,rectangle]{$M$}; 
						 \draw[dashed] (1,0) node[above right]{$1$}|-(0,1) node[left]{$1$};
					\end{tikzpicture}
			\end{center}
			\item Biểu diễn  miền nghiệm của mỗi bất phương trình $2x+3y-1<0$.
			\begin{itemize}
				 \item Vẽ đường thẳng $\Delta\colon 2x+3y-1=0$ đi qua hai điểm $\left(\dfrac{1}{2};0\right)$ và $\left(0;\dfrac{1}{3}\right)$.
				\item Xét điểm $O(0; 0)\notin \Delta$.  Ta có $2\cdot0+2\cdot0-1=-1<0$.
				\item Do đó, miền nghiệm của bất phương trình $2x+2y-1<0$ là nửa mặt phẳng bờ $\Delta$ chứa điểm $O$ không kể  đường thẳng $\Delta$.
			\end{itemize}
			\begin{center}
				\begin{tikzpicture}[scale=.9,>=stealth,font=\footnotesize,line join=round,line cap=round]
						 \def\a{-2/3} \def\b{1/3}
						\def\xmin{-2} \def\xmax{4} 
						\def\ymin{-3} \def\ymax{1.5} 
						\draw[->] (\xmin,0)--(\xmax+.2,0)  node [below]{$x$};
						\draw[->] (0,\ymin)--(0,\ymax+.2) node [left]{$y$}; 
						\fill (0,0) circle (1pt) node[below  left]{$O$};
						 \clip (\xmin,\ymin) rectangle (\xmax,\ymax);
						\fill[pattern=north east  lines] plot (\x,{\a*(\x)+\b})--(\xmax,\ymax) --cycle;
						\draw[smooth,samples=300] plot  (\x,{\a*(\x)+\b});
						\node[rotate=-35,right] at (.5,-.5) {$\Delta\colon 2x +3y -1=  0$};
						\draw[fill=black] (0,1/3) circle  (1pt) node[left]{$\tfrac{1}{3}$};
						\draw[fill=black] (0.5,0) circle (1pt)  node[below]{$\tfrac{1}{2}$};
					\end{tikzpicture}
			\end{center}
			\item Biểu diễn miền nghiệm  của mỗi bất phương trình $x+3y-2\le 0$.
			\begin{itemize}
				\item Vẽ đường  thẳng $\Delta\colon x+3y-2=0$ đi qua hai điểm $\left(2;0\right)$ và $\left(0;\dfrac{2}{3}\right)$.
				\item Xét  điểm $O(0; 0)\notin \Delta$. Ta có $0+3\cdot0-2=-2<0$.
				 \item Do đó, miền nghiệm của bất phương trình $x+3y-2\le 0$ là nửa mặt phẳng bờ $\Delta$ chứa điểm $O$ kể cả đường thẳng $\Delta$.
			\end{itemize}
			\begin{center}
				\begin{tikzpicture}[scale=.9,>=stealth,font=\footnotesize,line join=round,line cap=round]
						\def\a{-1/3}  \def\b{2/3}
						\def\xmin{-3.5} \def\xmax{4} 
						 \def\ymin{-2} \def\ymax{3}
						\draw[->]  (\xmin,0)--(\xmax+.2,0) node [below]{$x$};
						\draw[->] (0,\ymin)--(0,\ymax+.2)  node [left]{$y$};
						\fill (0,0) circle (1pt) node[below  left]{$O$};
						\clip (\xmin,\ymin) rectangle  (\xmax,\ymax);
						 \fill[pattern=north east lines] (\xmin,\ymax)--plot (\x,{\a*(\x)+\b})--(\xmax,\ymax) --cycle;
						\draw[smooth,samples=300] plot (\x,{\a*(\x)+\b}); 
						\node[rotate=-18,right] at (-3.5,1.5) {$\Delta\colon  x +3y -2= 0$};
						\draw[fill=black]  (0,2/3) circle (1pt) node[below right]{$\tfrac{2}{3}$};
						\draw[fill=black] (2,0) circle (1pt) node[below]{$2$}; 
					\end{tikzpicture}
			\end{center}
			\item Biểu diễn  miền nghiệm của mỗi bất phương trình $\dfrac{x+y}{2}>\dfrac{2x-y+1}{3}$.\\
			 Ta có
			\begin{eqnarray*}
				&&  \dfrac{x+y}{2}>\dfrac{2x-y+1}{3}\\
				&\Leftrightarrow& 3x+3y  >4x-2y+2\\
				&\Leftrightarrow& x-5y+2<0. 
			\end{eqnarray*}
			\begin{itemize}
				\item Vẽ đường  thẳng $\Delta\colon x-5y+2=0$ đi qua hai điểm $\left(-2;0\right)$ và $\left(0;\dfrac{2}{5}\right)$.
				\item Xét điểm $O(0; 0)\notin \Delta$. Ta có  $0-5\cdot0+2=2>0$.
				\item Do đó, miền nghiệm của bất phương trình $x-5y+2<0$ là  nửa mặt phẳng bờ $\Delta$ không chứa điểm $O$ không kể đường thẳng $\Delta$.
			\end{itemize}
			\begin{center}
				\begin{tikzpicture}[scale=.9,>=stealth,font=\footnotesize,line join=round,line cap=round]
						\def\a{1/5} \def\b{2/5} 
						 \def\xmin{-3.5} \def\xmax{4}
						\def\ymin{-2.5} \def\ymax{2} 
						\draw[->] (\xmin,0)--(\xmax+.2,0) node  [below]{$x$};
						 \draw[->] (0,\ymin)--(0,\ymax+.2) node [left]{$y$};
						\fill  (0,0) circle (1pt) node[below left]{$O$};
						\clip  (\xmin,\ymin) rectangle (\xmax,\ymax);
						 \fill[pattern=north east lines] plot (\x,{\a*(\x)+\b})--(\xmax,\ymin)--(\xmin,\ymin) --cycle;
						\draw[smooth,samples=300] plot  (\x,{\a*(\x)+\b});
						 \node[rotate=12,right] at (.5,.7) {$\Delta\colon x -5y +2= 0$};
						\draw[fill=black] (0,.4) circle (1pt) node[above left]{$\dfrac{2}{5}$}; 
						\draw[fill=black] (-2,0) circle (1pt)  node[above]{$-2$};
					\end{tikzpicture}
			\end{center}
			\item Biểu diễn miền nghiệm của mỗi bất phương trình $\dfrac{5x+y-1}{2}\ge  x+y+1$.\\
			 Ta có
			\begin{eqnarray*}
				&&  \dfrac{5x+y-1}{2}>x+y+1\\
				&\Leftrightarrow& 5x+y-1 >2x+2y+2\\ 
				&\Leftrightarrow&  3x-y-3\ge 0.
			\end{eqnarray*}
			\begin{itemize}
				\item Vẽ đường thẳng $\Delta\colon 3x-y-3=0$ đi qua hai điểm $\left(1;0\right)$  và $\left(0;-3\right)$.
				\item Xét điểm $O(0; 0)\notin \Delta$. Ta có  $3\cdot0-0-3=-3<0$.
				\item Do đó, miền nghiệm của bất phương trình $3x-y-3\ge 0$ là nửa mặt phẳng bờ $\Delta$  chứa điểm $O$ kể cả đường thẳng $\Delta$.
			\end{itemize}
			\begin{center}
				\begin{tikzpicture}[scale=.9,>=stealth,font=\footnotesize,line join=round,line cap=round]
						\def\a{3} \def\b{-3} 
						\def\xmin{-3} \def\xmax{3} 
						\def\ymin{-4} \def\ymax{2} 
						\draw[->] (\xmin,0)--(\xmax+.2,0)  node [below]{$x$};
						\draw[->] (0,\ymin)--(0,\ymax+.2)  node [left]{$y$};
						\fill (0,0) circle  (1pt) node[below left]{$O$};
						\clip  (\xmin,\ymin) rectangle (\xmax,\ymax);
						\fill[pattern=north east lines] (\xmin,\ymin)--plot  (\x,{\a*(\x)+\b})--(\xmin,\ymax) --cycle;
						\draw[smooth,samples=300] plot (\x,{\a*(\x)+\b}); 
						\node[rotate=70,right] at (.5,-3.2) {$\Delta\colon 3x  -y -3= 0$};
						\draw[fill=black] (0,-3) circle (1pt)  node[below right]{$-3$};
						\draw[fill=black] (1,0)  circle (1pt) node[above right]{$1$};
					\end{tikzpicture}
			\end{center}
		\end{enumerate}
	} 
\end{bt}

% nguồn: MaHoaNH_TuLuan_Lop10
\begin{bt}%[0D2V1-2]%[ma:508ea489]
	Cho tam giác $ABC$ có $A\left(1;2\right)$, $B\left(-3;-1\right)$ và $C\left(3;-4\right)$. Tìm điều kiện của tham số $m$ để điểm  $M\left(m;\dfrac{m-5}{3}\right)$ nằm bên trong tam giác $ABC$.
	\loigiai{ 
		\begin{center}
			\begin{tikzpicture}[scale=.8,>=stealth,font=\footnotesize,line join=round,line cap=round]
					 \path
						(1,2) coordinate (A) 
						 (-3,-1) coordinate (B)
						(3,-4) coordinate (C) 
						(1,-1)  coordinate (M);
					\draw[fill=cyan!50](A)--(B)--(C)--cycle; 
					\draw[->]  (-4,0)--(4,0) node[below]{$x$};
					\draw[->] (0,-5)--(0,3) node[left]{$y$}; 
					\node at  (0,0) [below left]{$O$};
					 \foreach \x/\g in {A/90,B/180,C/-45,M/0}
					\fill[black] (\x) circle (1pt) 
						 ($(\g:3mm)+(\x)$) node {$\x$};
				\end{tikzpicture}
		\end{center}
		Phương trình đường thẳng $AB\colon\dfrac{x-1}{-3-1}=\dfrac{y-2}{-1-2}\Leftrightarrow 3x-4y+5=0$.\\ 
		Phương trình đường thẳng  $BC\colon\dfrac{x+3}{3+3}=\dfrac{y+1}{-4+1}\Leftrightarrow x+2y+5=0$.\\
		Phương trình đường thẳng $AC\colon\dfrac{x-1}{3-1}=\dfrac{y-2}{-4-2}\Leftrightarrow  3x+y-5=0$.\\
		Điều kiện cần và đủ để điểm $M$ nằm bên trong tam giác $ABC$ là điểm $M$ cùng với mỗi đỉnh $A$, $B$, $C$ lần lượt  cùng phía với nhau đối với cạnh $BC$, $CA$, $AB$.
		\begin{itemize}
			 \item $M$ và $A$ nằm cùng phí đối với $BC$ $\Leftrightarrow\left(1+4+5\right)\left(m+2\dfrac{m-5}{3}+5\right)>0\Leftrightarrow m>-1$.
			\item $M$  và $B$ nằm cùng phí đối với $AC$ $\Leftrightarrow \left(-9-1-5\right)\left(3m+\dfrac{m-5}{3}-5\right)>0\Leftrightarrow m<2$.
			\item $M$ và $C$ nằm cùng phí đối với $AB$ $\Leftrightarrow \left(9+16+5\right)\left(3m-4\dfrac{m-5}{3}+5\right)>0\Leftrightarrow m>-7$. 
		\end{itemize}
		Suy ra $-1<m<2$ thỏa mãn điểm $M\left(m;\dfrac{m-5}{3}\right)$ nằm bên  trong tam giác $ABC$.
	 }
\end{bt}

% nguồn: MaHoaNH_TuLuan_Lop10
\begin{bt}%[0D2V1-2]%[ma:2fd5323b]
	Cho  tam giác $ABC$ có $A(0;3), B(-1;2), C(2;1)$. Tìm điều kiện của tham số $m$ để điểm $M\left(m;\dfrac{2m-1}{2} \right)$ nằm bên trong tam giác $ABC$.
	\loigiai{ 
		\begin{center}
			\begin{tikzpicture}[scale=1,>=stealth,font=\footnotesize,line join=round,line cap=round]
					 \path
						(0,3) coordinate (A) 
						(-1,2) coordinate  (B)
						(2,1)  coordinate (C)
						(0,2) coordinate (M); 
					\draw[fill=cyan!50](A)--(B)--(C)--cycle; 
					 \draw[->] (-2,0)--(3,0) node[below]{$x$};
					 \draw[->] (0,-1)--(0,4) node[left]{$y$};
					\node at (0,0) [below left]{$O$}; 
					\foreach \x/\g in {A/30,B/180,C/-45,M/0} 
					 \fill[black] (\x) circle (1pt)
						($(\g:3mm)+(\x)$) node  {$\x$};
				\end{tikzpicture}
		\end{center}
		Đường thẳng $AB\colon  \dfrac{x-0}{-1-0}=\dfrac{y-3}{2-3}\Leftrightarrow x-y+3=0$.\\
		 Đường thẳng $BC\colon \dfrac{x+1}{2+1}=\dfrac{y-2}{1-2}\Leftrightarrow x+3y-5=0$.\\
		Đường thẳng $AC\colon \dfrac{x-0}{2-0}=\dfrac{y-3}{1-3}\Leftrightarrow x+y-3=0$.\\ 
		Điều kiện cần và đủ để điểm $M$ nằm bên trong tam giác $ABC$ là điểm $M$ cùng với mỗi đỉnh $A$, $B$,  $C$ lần lượt cùng phía với nhau đối với cạnh $AB$, $AC$, $BC$.
		\begin{itemize}
			\item $M$ và $A$ nằm cùng phí đối  với $BC$
			 $\Leftrightarrow\left(0+3\cdot 3-5\right)
			\left(m+3\cdot \dfrac{2m-1}{2}-5\right)>0 
			 \Leftrightarrow m>\dfrac{13}{8}$.
			\item $M$  và $B$ nằm cùng phí đối với $AC$
			 $\Leftrightarrow \left(-1+2-3\right)
			\left(  m+\dfrac{2m-1}{2}-3\right)>0
			\Leftrightarrow m<\dfrac{7}{4}$. 
			\item $M$ và $C$ nằm cùng phí đối với $AB$ 
			$\Leftrightarrow  \left(2-1+3\right)
			 \left(m-\dfrac{2m-1}{2}+3\right)>0
			\Leftrightarrow  14>0$.
		\end{itemize}
		Suy ra $\dfrac{13}{8}<m<\dfrac{7}{4}$ thỏa mãn điểm  $M\left(m;\dfrac{2m-1}{2} \right)$ nằm bên trong tam giác $ABC$.
	 }
\end{bt}

% nguồn: MaHoaNH_TuLuan_Lop10
\begin{bt}%[0D2B1-2]%[ma:695512b7]
	Bình đến văn phòng phẩm mua quà tặng bạn. Trong cửa hàng có $3$ mặt hàng: bút, vở và thước. Bút có $5$ loại, vở có $4$ loại, thước có $3$ loại. Hỏi Bình có bao  nhiêu cách chọn một phần quà gồm $1$ bút, $1$ vở và $1$ thước?
	 \loigiai{
		Số cách để Bình chọn $1$ cây bút là $5$, chọn $1$ quyển vở là $4$ và $1$ cây  thước là $3$.\\
		Mỗi lần chọn một món đồ dùng là một công đoạn chọn quà của  Bình.\\
		Do đó số cách chọn $1$ phần quà gồm đủ $3$ món đồ đó là $5\cdot  4\cdot 3=60$ cách.
	} 
\end{bt}
